% before: pd-cm-lattice
The theta functions whose integrality drives this chapter live on elliptic curves with complex
multiplication by the ring of integers of an imaginary quadratic field. Such a curve is given
analytically by a lattice in $\mathbf C$ stable under that ring, and its invariants are the
Eisenstein series of the lattice.

% before: pd-j-invariant-integral
Integrality of Taylor coefficients at a prime $p$ needs a model of the curve with good reduction
at $p$. The first step towards such a model is that the $j$-invariant of a CM lattice is an
algebraic integer, so that the curve has potentially good reduction everywhere.

% before: pd-good-model
At a prime $l$ the lattice is moved into $\mathbf C_l$ by a field isomorphism, and a rescaling by
$\lambda$ changes the invariants by $\lambda^4$ and $\lambda^6$. A good model is a rescaling whose
invariants are integral over the maximal unramified extension and whose discriminant is a unit.

% before: pd-inert-formal-group
At a prime $l$ inert in the CM field the reduction is supersingular, and the formal group has
height two. Its logarithm can then be normalized to $\sum_jX^{l^{2j}}/l^j$, with denominator $l^j$
at the exponent $l^{2j}$. This shape underlies the sum over odd powers of $p$ in the gains at
inert primes (Definition~\ref{pd-class-sum}).

% before: pd-theta-function
The automorphic forms on definite unitary groups that are lifted along the chain are built from
theta functions in several variables attached to a Hermitian form and a self-dual lattice. In
one variable, the normalized theta function of a CM lattice is a modified sigma function, and it
is the building block of the comparison below.

% before: pd-cm-datum
The lattices of the host are compared with products of the one-variable CM lattice. The sum of
four copies of the Hermitian space is isometric to the standard space of the CM curve, by an
isometry carrying the fourth power of the lattice to a lattice commensurable with the standard
one. Local normalization asks that the fourth power of the theta function be, up to a $p$-integral
constant, the transfer of the product of one-variable theta functions through that isometry.

% before: pd-integral-jets
Divisibility of periods is measured on derivatives at the origin, which are the Taylor
coefficients multiplied by $\kappa!$. Integral jets ask for integrality of these derivatives, a
condition that allows the Taylor coefficients denominators up to $\kappa!$, and the jet pairing
evaluates a differential operator with integral coefficients at the origin.

% before: pd-jets-closure
The proof works with series whose jets are integral, and it combines them by sums, products,
exponentials and linear changes of variables. These operations preserve integral jets, and a series
whose fourth tensor power has integral jets has integral jets itself, which is how the fourfold
products of the local normalization are undone.

% before: pd-normalizing-series
The normalized theta function of a CM lattice is a translate of $\sigma^*_L$, and in the coordinate
of the formal group $\sigma^*_L$ is $X$ times $e^{-\phi_L}$. The integrality of the theta jets
comes down to the integrality of the jets of $\phi_L$, measured with the invariant derivation.

% before: pd-cm-isogeny
The jets of $\phi_L$ are controlled by an endomorphism of the curve whose degree is prime to $p$.
Its kernel gives a distribution relation for $\wp_L$, and comparing quasi-periods turns that relation
into an identity for the constant $s_2(L)$ that enters $\sigma^*_L$.

% before: pd-cm-isogeny-classical
The identity for $s_2(L)$ is classical. The textbook route integrates the distribution relation and
fixes the constant of integration, which the proof above avoids.

% before: pd-sigma-congruences
In the coordinate $X$ these identities become congruences between series whose differences have
integral exponentials. The quadratic term with $s_2(L)$ cancels in each of them, because the identity
for $s_2(L)$ was built into $\sigma^*_L$. Additivity in the first variable comes from the addition
formula of $\sigma_L$ by a halving argument, which needs $|2|=1$.

% before: pd-jets-abstract
The three congruences determine the jets. Parity controls the first jet, the endomorphism
controls the second through the unit $a-b$, and additivity of $\beta$ in each variable makes the
third invariant derivative of $\phi$ integral, which gives all the higher jets.

% before: pd-cm-jets
For CM by $\mathcal O_K$ at an odd prime not dividing $d$, the endomorphism $m+\omega_K$ with
$p=2m+1$ has degree $(p^2+d)/4$, prime to $p$, and $\alpha-\bar\alpha=i\sqrt d$ is a unit. So the
abstract statement applies to the normalizing series of the CM lattice.

% before: pd-translation-system
The theta function of the statement is compared with translates of the one-variable theta
function at points of order prime to $p$. Divided by $e^{-\phi_L}$, these translates satisfy the
addition formula of $\sigma_L$ written in the formal group coordinate, and that formula is all the
integrality argument uses.

% before: pd-translation-units
At a point of order $n$ prime to $p$, iterating the addition formula $n$ times returns to the
origin, where the series is a unit multiple of $X$. Comparing lowest terms then forces the constant
term to be a unit and the normalized series to be an integral $n$-th root of an integral series.

% before: pd-analytic-translation-system
The translates of $\sigma^*_L$ satisfy the axioms of a translation system for the analytic
parametrization of $\mathcal E_L$. Torsion points of order prime to $p$ have integral coordinates,
and the operators $U^L$ compose up to roots of unity, which lets a translate of
$\vartheta_{L,\omega_0/2}$ be written as a translate of $\sigma^*_L$.

% before: pd-theta-translates
The two preceding lemmas give the integrality of every translate at a point of order prime to $p$.
The point $\omega_0/2$ has order $2$, so the normalization of $\vartheta_{L,\omega_0/2}$ is a unit
when $p$ is odd.

% before: pd-fourfold-product
Local normalization writes the fourth tensor power of a translate of $\theta$ as a combination of
products of one-variable translates. Its coefficients are units or roots of unity at $p$, so the
fourfold product inherits the shape of the one-variable translates.

% before: pd-multiplier-coordinates
The variables of the statement enter through vectors of the multiplier module. Their coordinates in
the basis given by the isometry of the local normalization are $p$-integral, because the index of
condition (2) is prime to $p$ and the multiplier ring of $L$ is $\mathcal O_K$.

% before: pd-jets-from-fourth-power
The last step changes variables in the fourfold product and takes a fourth root. Integrality of
the jets of $\phi_L$ gives integral jets in the logarithmic coordinates, and the integral matrix
of the previous lemma carries them to the variables of the statement.

% before: pd-theta-jets-integral
With these lemmas the main local statement follows: restricted to linear combinations of
vectors in the multiplier module, combinations of $p$-integral torsion translates of a locally
normalized theta function have integral jets at every odd prime of good reduction not dividing the
discriminant. This is the integrality that the kernels of the lifts inherit at the served places.

% before: pd-period-data
The periods of the iterated lifts are computed from the data of a definite-side datum: the
coefficient field and its places, the stabilizers of the lattice classes, the top forms and the
point lattices on which the lifted forms are evaluated. The statements of this chapter use these
parts of the datum and refer to Chapter~\ref{chapter-unitary-theta} for the rest.

% before: pd-period-integrality
Each lift averages over lattice classes with weights the inverse orders of their stabilizers, so
periods of integral vectors are integral only up to the stabilizer denominator. Period
integrality bounds the denominator of the period of a point-lattice vector against an integral top
form by that denominator, at every finite place.

% before: pd-twisted-period-integrality
The integral structure of Chapter~\ref{chapter-integral-structure} glues the lifts across tuples
of Hecke characters, so the twisted periods need the same control. The condition is imposed only
at places whose residue characteristic is not exceptional for the tuple; the multiplicity bound
uses it at a place $w$ only for the tuple $\mu_w$, for which $p_w$ is never exceptional.

% before: pd-size-data
The multiplicity bound compares valuations of periods with numerical quantities fixed in advance:
the gains of the lifts at each place, the local terms at the twisted places and the local errors.
These belong to the size data of the datum, which also carry the scale, the exponents $a_t$ and the
weights used at inert primes.

% before: pd-multiplicity-bound
In Stepanov's method an auxiliary polynomial vanishes to high order at the points it is built to
vanish on. Here the corresponding statement is that every period of a lifted division vector is
divisible at each place by the accumulated gain of the $m$ lifts, corrected by the local terms.

% before: pd-lift-valuation
The divisibility of a period is obtained by first bounding the valuation of the iterated lift
itself, as a vector of the top lattice. The following notion records that a lifted vector is a
multiple of a locally integral vector by a scalar of controlled valuation.

% before: pd-period-of-divisible-lift
The next lemma converts a valuation bound for the iterated lift into a valuation bound for its
periods. It loses only the stabilizer denominator, by period integrality, and is the common step in
the three cases of the proof of the multiplicity bound.

% before: pd-family
The multiplicity bound, the gain link and the links of later chapters are proved for all large
scales of suitable families of definite-side data, uniformly in the members at each scale. Their
hypotheses are statements about families that hold once the constants of the family are large
enough, and largeness is arranged in stages because some constants must be chosen after others.

% before: pd-hecke-at-T
At the auxiliary prime $T_u$ the iterated lift is bounded on a lattice obtained from the
ramified forms by Hecke operators at the places attached to top embeddings. The exponent $E_T$
collects the distances of the weights from the middle of each level, and it is the local error
at $T_u$ in the multiplicity bound.

% before: pd-orbit-places
Away from $T_u$ and the level primes the bounds on the lifts come from the expansions of the
theta kernels at each place. At a prime $p$ the places of the datum are permuted by translation by
the Frobenius element $h_p$, and the shape of each lift at a place is a Young diagram read off from
the weights. The local theory works one orbit of this translation at a time.

% before: pd-minor-valuation
The lattices of forms at a free place are cut out by products of minors whose shape is the Young
diagram of the row exponents. A lift therefore gains valuation $B$ when its theta kernel, paired as
a differential operator with these products of minors, gives values divisible by $l^{B}$.

% before: pd-kernel-entries
The kernel entries make that reading of the lattices of forms precise. In coordinates on the source
and target lattices, each kernel map becomes the pairing of its theta kernel series with products
of minors, up to a constant.

% before: pd-orbit-expansions
At an inert prime the theta kernel, written in the Lubin-Tate logarithm of the prime, splits along
the Frobenius orbits into an integral factor and an exponential. The conditions (B1) to (B4) say
that the exponent is defined over the unramified extension, invariant under the torus, of high
enough bidegree, and nearly additive in each parity.

% before: pd-local-model
At an inert prime $p$ the lattices of forms of a word are compared with lattices of $p$-adic forms
on a product of unitary groups that split over unramified extensions of $\mathbf Q_p$, one for each
Frobenius orbit of the embeddings. The local model records these lattices, their level groups and
the Hecke operators at $p$, and its fills give the weights of the successive lifts.

% before: pd-lattice-instance
The local models are attached to the definite-side datum at the admissible places, and their
positions are identified with the embeddings so that each Frobenius orbit of the gain data is one
orbit of the model. The comparison carries the lattice of forms in degree $0$ onto the integral
module of the model, so that bounds proved in the model transfer to the iterated lift.

% before: pd-formal-group-system
In the local model the lifts are given by kernels whose expansions in the Lubin-Tate logarithm
split into an integral factor and the exponential of a bi-additive series. A formal group system
records these kernels with the representatives of the lattice classes and the conditions that let
the jets of the kernels be bounded orbit by orbit.

% before: pd-served-selection
A served selection chooses kernel entries whose theta kernels have these expansions at every free
place, together with the orbit bounds that the local theory of Frobenius orbits provides. Among the
conditions is the jet integrality of Theorem~\ref{pd-theta-jets-integral}.

% before: pd-kernel-form
A bound on the theta kernels against products of minors becomes a bound on the lift through the
coordinates of the kernel entries. The only loss is the stabilizer denominator, which comes from the
weights $1/|\Gamma_c|$ in the lift and disappears when the stabilizers are trivial.

% before: pd-odd-level-programme
The weight of a lift at the positions of an orbit can be spent at the odd distances along the orbit,
at cost $l^i$ at the source and $1$ at the target. The resulting linear programme measures how much
of the gain $\gamma_l$ per unit of weight is lost to these exchanges.

% before: pd-lp-duality
The comparisons of the programme for different weights go through its dual. An optimal dual
vector prices each position, and the positions with a high price carry little weight.

% before: pd-lp-comparison
In the application the weights at the positions are close to $(k-1)$ times a fixed vector, with a
defect of $2w$ at the top positions and an error of at most $4w$. The programme for these weights
loses at most $11/l$ of the total weight against the programme for $(k-1)w$.

% before: pd-weights
The bound for an orbit expansion is a Dwork-type estimate: each coefficient of a power series is
compared with a power of $|l|$ that depends linearly on the exponent. Assigning a weight to each
variable keeps track of these powers through substitutions, products and exponentials.

% before: pd-weight-calculus
These rules are the whole calculus of weights that the estimate uses. The factorials enter because
the jet pairing multiplies each coefficient by $\kappa!$, and the exponential and the logarithm move
between the series $\beta$ and its exponential.

% before: pd-letter-series
The conditions (B1) to (B4) on the series $\beta$ of an orbit expansion are used in a form that
does not refer to the theta kernel. A bipower is a monomial at which the coefficient of $\beta$ may
fail the better of the two weight bounds below.

% before: pd-log-inverse
The orbit expansion is written with the orbit logarithm substituted, so the estimate needs its
inverse. The inverse of a logarithm of height two has controlled denominators, of size $\rho_l$ per
unit of degree.

% before: pd-binomial-extraction
The exponentials of $\delta_{\mathrm{ev}}\beta$ and $\delta_{\mathrm{od}}\beta$ are integral by
(B4), but the estimate needs bounds on $\beta$ itself. The binomial coefficients that appear in
these differences recover the coefficients of $\beta$, with one unit of loss at the bipowers.

% before: pd-letter-weights
The weights of a letter series follow from its integrality conditions. The torus invariance and the
unramified coefficients quantize the valuations, which turns the bound $\rho_l$ from the inverse
logarithm into the bound $1/(l-1)$.

% before: pd-letter-substitution
In the theta kernel the series $\beta$ is evaluated at linear forms in many variables. The loss of
one unit survives only at the products of two pure $l$-power monomials, and in a frame only at
positions of different parity whose rows are not nearly orthogonal.

% before: pd-weighted-dwork
Weights on the variables can absorb the loss of one unit at the remaining monomials. Each unit of
weight on a variable then costs one unit of the bound per unit of its exponent, against the gain
$\gamma_l$ from the factorials.

% before: pd-one-place-weighted
At a split orbit there is one place and every pair of exponents is a bipower. A constant weight
$1/(l+1)$ absorbs the loss there and leaves the gain $\rho_l$ per unit of degree.

% before: pd-dual-radii
The weights for an orbit of even length are built from a dual vector of the odd-level programme and
from numbers at the top positions. These inequalities are the cases of the hypothesis of
Proposition~\ref{pd-weighted-dwork} that they must satisfy.

% before: pd-row-radii
The numbers at the top positions are chosen one position at a time, against the costs of the labels
that point to them. The resulting bound is the sum of the terms $\operatorname{top}(s)$, the corrections
$\gamma_lW(s)$ and the linear programme $F(D)$ that appear in the orbit bound.

% before: pd-single-monomial
A frame whose top rows are nearly orthogonal to the rows that do not carry labels makes every
loss of one unit occur at a labelled pair. The row radii then give the bound for a single monomial,
which is the estimate behind the orbit bound.

% before: pd-orbit-letter-bounds
For a single Frobenius orbit of even length, an orbit expansion bounds the theta kernel below by
the terms $a/p+a/p^3+\cdots$ of the announcement, corrected by the caps of the neighbouring
positions and by a linear programme for the remaining weight. The split orbits give the smaller
gain $|\nu|/(p^2-1)$.

% before: pd-served-lattice
At a prime of inert class with served representatives, the served column bound (ii) of
Theorem~\ref{pd-multiplicity-from-integrality} exceeds the local gains. At each served orbit it
gains the unit gain, in which the top classes at even positions of the orbit contribute without
caps, less the loss of a Hecke operator. This gain holds on a sublattice of the forms, the served
local lattice, built from products of Hecke operators at the served orbits; an orbit left out of a
product is paid for by a power of $p$.

% before: pd-formal-group-numbers
At a served orbit the bound comes from a formal group system on the local model of the place. It
is stated through three numbers of one orbit of the system, computed from the weights of the fills:
the gains of the odd steps along the orbit, with and without the steps that land on top positions
of the other parity, and the loss of the Hecke operator of the orbit. The digit error accounts for
the $p$-adic digits of these weights.

% before: pd-hecke-chain-operators
The proof of Theorem~\ref{pd-formal-group-bounds} moves a power of the Hecke operator $y_v$ at level $0$
through the chain to an operator at level $m_*$ built from the edge operators of the orbit. That operator
is a combination of edge operators of indices $0$ to $r$, with coefficients given by complete homogeneous
symmetric polynomials in elements of known valuation. Its terms carry the powers of $p$ that appear in
condition~(i).

% before: pd-kernel-orbit-index
Condition~(i) attaches to an edge of index $i$ at an orbit a power of $p$ on each minor at the positions
of that orbit. The kernel series is bounded one orbit at a time against these weighted minors. With index
$0$ these bounds are the lower bounds of Definition~\ref{pd-formal-group-system}.

% before: pd-kernel-orbit-series
At one orbit the kernel series is controlled through a single series $H_u$ in the orbit coordinates. A
series in the class of $H_u$ is a sum of powers of $H_u$ divided by the same powers of $p$, with integral
cofactors. After the substitution $Z\mapsto\pi Z$ with $\pi^{p^2-1}=p$, the kernel series at an orbit lies
in such a class.

% before: pd-hecke-chain-edge-scaling
An edge operator replaces a lattice at one orbit by a sublattice, and the lattices $H^\Delta$ of the new
tuple are larger. A power of $p$ fixed by the weights brings them back, because the weights bound the
degree of the minors in the rows that the edge divides by $p$. So the edge operators preserve the span of
the integral module, on which the expansion of the chain holds.

% before: pd-hecke-chain-edge-count
Moving an edge operator past a map of the chain amounts to counting lattices. A matrix carries a lattice
into a sublattice of index $q^i$ of a second lattice, or a superlattice of index $q^{i-l}$ of the first
lattice into the second; the two counts differ by Gaussian binomial coefficients.

% before: pd-hecke-chain-relation
The Weil index formula of condition~(f) makes the kernel depend on the lattices at an orbit only through
an indicator function and explicit powers of $p$, $u$ and $u'$. With the count of
Lemma~\ref{pd-hecke-chain-edge-count} and the adjunction of condition~(h), the edge operator of index $i$
at level $j+1$ after $K_j$ becomes $K_j$ after a combination of edge operators of index at most $i$ at
level $j$. The coefficients are elementary symmetric polynomials in $t+1$ elements.

% before: pd-hecke-chain-inversion
The relation of Lemma~\ref{pd-hecke-chain-relation} multiplies a generating series of operators by a
product of linear factors. Iterating it along the chain, and dividing by the product of these factors,
moves an operator at level $0$ to level $j$. The division uses complete homogeneous symmetric
polynomials.

% before: pd-hecke-chain-commutation
Combining the relation with its inversion at the index $r$ of the Hecke operators commutes the chain with
$y_v$ up to units. The coefficients of the resulting orbit operator have valuation at least $0$, and the
coefficient of its term of index $0$ has valuation at least the even top sum, which the bound of index $0$
in Lemma~\ref{pd-kernel-orbit-bounds} loses.

% before: pd-kernel-orbit-descent
The bounds on the kernel hold over the field $\mathcal K_c$, while the integral modules are defined over
$\mathcal K$. A polynomial over $\mathcal K$ that is a combination of elements of $H^\Delta(L)$ with
coefficients of valuation at least $g$ in $\mathcal K_c$ is such a combination with coefficients in
$\mathcal K$.

% before: pd-kernel-chain-bound
The expansion of condition~(e) writes each map of the chain as a sum of kernels with the orders of
stabilizers as denominators, and condition~(i) controls these kernels after the weighted edge operators.
A lower bound on the pairings of the weighted minors with the kernel series therefore bounds the chain
followed by these operators, at a cost of $d_\Gamma(j)$.

% before: pd-kernel-orbit-product
Condition~(b) splits every kernel series, up to an integral factor, into a product over the orbits of
exponentials of the bi-additive series. The jet pairing of a product of minors with such a product is a
product over the orbits, so rational lower bounds at the orbits add.

% before: pd-orbit-class-law
At one orbit the kernel series involves the exponential of a sum of values of the bi-additive series
$\beta$, and Lemma~\ref{pd-kernel-orbit-class} needs this exponential in closed form. After the
logarithm $\ell$ is substituted, additivity becomes additivity for the formal group law with logarithm
$\ell$, and $\beta$ is compared with series built from $\ell$ and the odd logarithm $\varrho$. The
comparison is made up to integral exponentials, through coboundaries and a Frobenius twist.

% before: pd-orbit-class-facts
The comparison rests on standard properties of the formal group law and of the two logarithms. Dwork's
lemma turns the integrality of an exponential into a congruence modulo $p$ for the Frobenius twist, and a
divided derivative bound controls a series when one of its arguments moves by a multiple of $p$.

% before: pd-orbit-class-one-variable
A series in one variable that is additive for $F$ up to an integral series is determined, up to an
integral series, by two constants. Its coefficients are integral except at the powers of $p$, where the
binomial coefficients force a correction by $\varrho$ or by $\ell$.

% before: pd-orbit-class-two-variables
In two variables the same rigidity holds: a series vanishing on the axes and additive in each variable up
to integral series is a multiple of $\ell(X_0)\ell(X_1)$ up to an integral series. Slicing in either
variable reduces it to Lemma~\ref{pd-orbit-class-one-variable}, and the two slicings constrain each other.

% before: pd-orbit-class-average
The comparison of $\beta$ with an explicit series starts with an average over the roots of unity of order
$p^2-1$. The average keeps the relations $\simeq_0$, as $p^2-1$ is prime to $p$, and the twisted
invariance under these roots of unity removes the monomials whose exponents break the congruence
$n_0\equiv pn_1$.

% before: pd-orbit-class-derivative
Differentiating the coboundary relations in one variable at a point where they vanish gives a series in
one variable that is additive up to an integral series. Lemma~\ref{pd-orbit-class-one-variable} writes it
through $\varrho$ and $\ell$, and the balanced support leaves only the multiple of $\varrho$.

% before: pd-orbit-class-frobenius-twist
The Frobenius twist $\mathrm D_\sigma$ is the operator of Dwork's lemma, so the relation
$\delta_FG\simeq_00$ gives a congruence for the coboundary of $\mathrm D_\sigma G$. Its defect is
$\sigma(c)$ times the coboundary of $\varrho$, with $c$ the derivative of $G$.

% before: pd-orbit-class-integral
With these defects, the coboundaries of $\mathrm D_\sigma G$ and of its exchange are integral up to
explicit multiples of $\delta_F\varrho$. Subtracting the matching multiples of $\varrho(X_0)\ell(X_1)$
and $\ell(X_0)\varrho(X_1)$ gives a series to which Lemma~\ref{pd-orbit-class-two-variables} applies, and
the balanced support excludes the term $\ell(X_0)\ell(X_1)$.

% before: pd-orbit-class-constants
The integrality of the corrected twist relates the derivatives of $G$ and of its exchange. The
coefficients of $X_0X_1^{p^{2i+1}}$ in the two integral series bound a fixed combination of $c$ and
$\sigma(c')$ times $p^{-(i+1)}$ for every $i$, so the combination vanishes.

% before: pd-orbit-class-comparison
When $\sigma^2(c)=c$, the corrected twist of $G$ is the Dwork defect of $G-G_{c,\sigma(c)}$, because the
explicit series has the same twist up to the two correction terms. Integrality of that defect is then
the relation $G\simeq_0G_{c,\sigma(c)}$.

% before: pd-orbit-class-classification
These lemmas combine into a classification: a series additive for $F$ in each variable up to $\simeq_0$,
and twisted by the roots of unity of order $p^2-1$, is equivalent to an explicit series
$G_{c,\sigma(c)}$. The constant $c$ is the derivative of the averaged series, and $\sigma^2(c)=c$ comes
from comparing the series with its exchange.

% before: pd-orbit-class-beta
A bi-additive series becomes, after the logarithm $\ell$ is substituted, a series satisfying the
hypotheses of the classification, because $\ell$ turns the formal group law into addition. This gives
the constant $c$ of part~(1) of Lemma~\ref{pd-kernel-orbit-class}.

% before: pd-orbit-class-series
The exponential of the explicit series is computed after the substitution $Z\mapsto\pi Z$. In the orbit
coordinates this substitution amounts to substituting $\pi\mu_\pi(L)$, with $\mu_\pi$ the inverse of the
rescaled logarithm, into the series in the logarithm, and congruences modulo $pR_\pi$ organize the
computation.

% before: pd-orbit-class-inverse
The rescaled logarithm has integer coefficients and reduces to $X+X^{p^2}$ modulo $p$. Its inverse is
therefore the alternating series modulo $p$, which produces the series $S_{s,a,b}$ of the definition of
$H_u$.

% before: pd-orbit-class-frobenius
Applied to a linear form in the variables of one parity, the inverse $\mu_\pi$ gives the series $S_t$
with Witt vector coefficients modulo $p$. The $p$-th power of such a form is computed modulo $p$ through
the Frobenius of the Witt vectors.

% before: pd-orbit-class-dual
The orbit coordinates are built from the elements $\omega$ and $\omega^*$ of the formal group system,
which are dual for the sum over the even powers of $\sigma$. The duality also holds in the other order, at
any two positions of the same parity.

% before: pd-orbit-class-two-sums
The leading terms of the explicit series, at the substituted orbit coordinates, sum to $H_u$ modulo $p$.
The duality of $\omega$ and $\omega^*$ collapses the double sum over positions to the products
$S_s^pS_{s+1}$ of consecutive positions.

% before: pd-orbit-class-leading
After the substitution of $\pi$, only the leading terms of the explicit series keep a denominator. Every
other term is $\pi^{p+1}$ times an integral one, because its degree grows faster than its denominator.

% before: pd-orbit-class-explicit
The leading terms and their sum modulo $p$ give the exponential of the explicit series at the substituted
orbit coordinates. The remainder is $\pi^{p+1}$ times a series with coefficients in $R_\pi$, and its
exponential has coefficients in $R_\pi$.

% before: pd-kernel-orbit-class
The exponential of the bi-additive series at one orbit is computed in the orbit coordinates after the
substitution $Z\mapsto\pi Z$. It is the exponential of $-(\pi^{p+1}/p)H_u$ times an integral series, and
the numbers $\pi^{(p+1)i}/i!$ are integral, so the kernel series at the orbit lies in the class of $H_u$.

% before: pd-orbit-generic-data
The bounds of Proposition~\ref{pd-kernel-orbit-generic} hold for any family of columns at the positions of
a cycle of even length, with any set of top positions. They depend on the columns only through three
counts at each position: the columns of height one at top positions, the other columns, and the cells of
the other columns. In these counts the gains take the form of the orbit gains of
Definition~\ref{pd-formal-group-numbers}.

% before: pd-orbit-generic-entry
At a generic frame a column of height one at a top position gives a linear form, and the other columns
stay minors. A product of powers of such linear forms with minors is an entry polynomial. Its Hasse
derivatives along column directions describe how it changes when multiples of the top rows are added
along single columns.

% before: pd-orbit-generic-hasse
Under the jet pairing, multiplying a series by a linear form is adjoint to differentiating the polynomial
along that form. This turns the linear factors of the decomposed series below into Hasse derivatives of
the minors, at the cost of factorials. The same adjunction for linear substitutions moves a change of
variables from the series to the polynomial.

% before: pd-orbit-generic-minors
A minor is linear in each of its columns, and adding multiples of one row vector along several columns
changes its matrix by a matrix of rank one. So a Hasse derivative of a product of minors along column
directions vanishes once it differentiates too often at one position, and the number of minors at that
position bounds the order.

% before: pd-orbit-generic-decomposition
Splitting off the top rows as new variables turns $H_u$, in bounded degree and modulo $p$, into monomials
in the new variables, products of such monomials with linear forms in the old variables, and a remainder
of high degree in the old variables. The congruence $(f+g)^p\equiv f^p+g^p$ modulo $p$ separates the two
kinds of variables in the powers of the series $S_{s,a,b}$.

% before: pd-orbit-generic-term
A series in the class of the decomposed $H_u$ is a sum of powers of the decomposition divided by powers of
$p$. Pairing one term of this expansion with a monomial in the new variables times a polynomial in the old
ones produces factorials, and the linear forms become Hasse derivatives. By the ultrametric inequality
one term bounds the whole pairing.

% before: pd-orbit-generic-factorials
Legendre's formula converts the factorials of a single term into powers of $p$. Their total, less the
powers of $p$ lost to the class, is at least $\gamma_p$ for each top cell and $\rho_p$ for each cell,
minus the exponents of the term and the digit error, which collects the losses from the digit sums.

% before: pd-orbit-generic-star
The decomposition, the bound for one term and the count of factorials combine into a bound for the pairing
of an entry polynomial with a series in the class of $H_u$. The bound still contains the exponents of the
chosen term, and the lemma also gives the constraints on them that the next lemmas use: they fit under the
exponents of the linear forms, and the Hasse derivative of the minors is nonzero.

% before: pd-orbit-generic-gain-sums
The exponents of a term are distributed over the top positions and the odd steps $2e+1$ along the cycle.
At each step they are bounded by the share $A(s)/p^{2e+1}$ of the top cells, and those of the Hasse
derivatives also by the number of columns at the end of the step. Summing these bounds over the steps gives
the excess gain and the capped gain.

% before: pd-orbit-generic-entry-bounds
Each exponent of a term is charged to one top position and one odd step. The constraints of
Lemma~\ref{pd-orbit-generic-star} and the vanishing of Hasse derivatives of minors then give the
hypotheses of Lemma~\ref{pd-orbit-generic-gain-sums}, in three ways that match the three bounds of
Proposition~\ref{pd-kernel-orbit-generic}.

% before: pd-orbit-generic-reduction
A coefficient of a frame jet is at most the value of the frame jet at a generic point, where the variables
of the frame are taken in a field of rational functions with the Gauss norm. At that point the product of
minors is an entry polynomial, whose linear forms come from the columns of height one at top positions.

% before: pd-orbit-generic-bounds
The reduction to a generic point and the bounds for entry polynomials give the bounds for the frame jets
of arbitrary column data. The substitution $Z\mapsto\pi Z$ multiplies a frame jet by $\pi$ to the number
of cells, which cancels the term $\rho_p$ for each cell in those bounds.

% before: pd-kernel-orbit-generic
For a series in the class of $H_u$, the pairings with the minors at one orbit are bounded by the gains of
the odd steps along the orbit. The three bounds serve an orbit outside a product of Hecke operators, an
orbit of the product in the term of index $0$, and an orbit of the product in a term of positive index.

% before: pd-kernel-orbit-bounds
Lemma~\ref{pd-kernel-orbit-class} and Proposition~\ref{pd-kernel-orbit-generic} combine into lower bounds
of the system at each orbit. They are joined through an absolute value on $\mathcal K_c$, extended to an
algebraic closure in which $p$ has a $(p^2-1)$-th root.

% before: pd-hecke-chain-step
One map of the chain is bounded after a product of orbit operators. Expanding the product into weighted
edge operators, each term is bounded by Lemmas~\ref{pd-kernel-orbit-product}
and~\ref{pd-kernel-chain-bound}, with the bounds of Lemma~\ref{pd-kernel-orbit-bounds} at the orbits of
the product and the given lower bounds at the others. At an orbit of the product, the coefficient of the
term of index $0$ makes up for the even top sum.

% before: pd-formal-group-bounds
An admitted formal group system bounds the chain of its local model below at every orbit. Products
of Hecke operators at a set of orbits raise this bound by the orbit gain less the orbit loss at each
orbit of the set, provided the other orbits keep their lower bounds. An orbit of the set that is
left out of a product is paid for by a power of the uniformizer.

% before: pd-served-numerics
Once a tie identifies the positions of an orbit of the local model with the Frobenius orbit of a
representative, the numbers of the gain data at the representative and those of the formal group
system at the orbit are made from the same weights. They compare in the direction that the formal
group bounds need, and for primes larger than $2a$ the gains of the gain data are not positive.

% before: pd-served-unit-bound
Condition~(1)(b) of Definition~\ref{pd-served-selection} provides, at each admissible place, an
admitted formal group system for the chain of a used tuple, with lower bounds at the orbits that
carry no representative. Theorem~\ref{pd-formal-group-bounds}, with the comparison of
Lemma~\ref{pd-served-numerics}, then bounds the chain applied to a product of Hecke operators at
representatives. The orbits without representatives contribute the base exponent.

% before: pd-served-lattice-bound
On the served local lattice, the power of $p$ attached to the orbits left out of a product makes up
for their missing unit gains. So every function of the lattice reaches the gain of each lift at
$p$.

% before: pd-served-law
Where no gain is available, the kernel bounds of the levels control the iterated lift. At signal
primes and primes of inert class the orbit bounds add the local gains, and on the served local
lattice the unit bound adds the served gain; the comparison of the twisted lattices carries the
bound for the local chain to the iterated lift.

% before: pd-multiplicity-from-integrality
The multiplicity bound is proved at all large scales of the designed families of
Chapter~\ref{chapter-assembly}. The places of the coefficient field fall into three kinds: those
over the auxiliary prime $T_u$, those over the level primes $g_0,g_1$, and the rest. At each kind
a different input bounds the valuation of the iterated lift, and period integrality turns the
bound into divisibility of the periods. Hypotheses (3) and (4), period integrality and its twisted
form, are proved for the families of the construction in
Section~\ref{section-pd-period-integrality}.

% before: pd-weight-lattices
The period of a point-lattice vector is a sum over point classes of pairings between vectors of a
polynomial representation of $GL_n$. Whether such a pairing is integral is decided by two lattices
in the representation attached to a lattice of $E^n$: the dual Weyl module lattice, which contains
the rows of integral top forms, and the Weyl module lattice of the dual lattice, which contains the
values of point-lattice vectors.

% before: pd-weight-duality
The two lattices are dual to each other under the normalized Fock pairing. This makes the pairing
of a point-lattice vector with an integral row an algebraic integer, place by place.

% before: pd-point-data
For the data of the construction the point classes are classes of lattices of rank $N(m)-1$, and
the period of the iterated lift is computed class by class. Each class contributes the pairing of
the value of the vector at the class with the row of the top form there, weighted by the inverse
order of the stabilizer of the class.

% before: pd-class-term-integral
The weight $|\Gamma_c|^{-1}$ is the only denominator in the term of a class. Without it, the term is
the pairing of a vector of the Weyl module lattice of $L_c^\vee$ with a vector of the dual Weyl
module lattice of $L_c$, which is integral by the duality.

% before: pd-true-columns
The members of the families of the construction are definite-side data assembled from explicit
classes of lattices, the true columns. For these data the period pairings of
Definition~\ref{pd-period-data} are the periods of point data, and period integrality reduces to
the integrality of the rows of the top forms.

% before: pd-standard-comparison
A datum built from true columns takes its period pairings from the point data of the tuples.
Conjugating the rows of a top form turns the period into a Hermitian product with a vector fixed by
the stabilizers of the point classes, which is a sum of one term for each class. This uses that the
stabilizers act through unitary matrices and characters of finite order. For a twisted tuple the
point classes carry frames at the twisting places of the true columns, and a tie tied to the true
columns ramifies over the same places.

% before: pd-period-integrality-from-rows
Summing the terms of the classes gives the period. The stabilizer of each point class injects into
the stabilizer of a lattice class of the top level, whose order divides the stabilizer denominator,
so the sum has the denominator allowed by period integrality. The same computation applies to the
twisted chains, whose stabilizer denominators agree with those of the base tuple once they are
defined.

% before: pd-comparison-class-split
The proof splits the period through the Hermitian product because that product is where the
comparison places the point classes. The period of point data is itself a sum over the point
classes, so the same class terms also come from its linearity alone.

% before: pd-row-stage
The top forms of the construction come from algebraic modular forms on lattices of rank $N(m)-1$
with values in $V_\sigma(L)$. Their integral structure consists of the forms whose value at every
lattice lies in the dual Weyl module lattice of that lattice, and the rows of a top form are its
values at the representatives of the point classes.

% before: pd-rows-integral-built
At a row stage the integral top forms come from $\nabla$-integral algebraic modular forms, so
their rows are integral. With Theorem~\ref{pd-period-integrality-from-rows} this gives period
integrality at every datum built at a row stage.

% before: pd-period-integrality-along-families
The multiplicity bound uses period integrality as a property of families at all large scales. A
slot over true columns places every member of every family built from it at a datum built from
true columns with integral period rows, so the integrality holds there at every scale.

% before: pd-integrality-inputs
The slots of the construction that the main argument uses lie over true columns, because their
cores are built at row stages. Theorem~\ref{pd-period-integrality-along-families} then supplies
hypotheses (3) and (4) of Theorem~\ref{pd-multiplicity-from-integrality} through the collection of
definite-side hypotheses of Chapter~\ref{chapter-assembly}.

% before: pd-class-sum
At a prime $p$ inert in the CM field the gain of a lift at a top class $t$ is compared with a sum
over odd powers of $p$, reflecting the height-two formal group. Only the odd powers $p^i\leq pa_t$
contribute, and a weight is subtracted at each of them.

% before: pd-gain-data
The gains at an inert prime $p$ are counted orbit by orbit, along the orbits of translation by
the $A$-component of the Frobenius class of $p$ on the top classes. The gain data choose
representatives of these orbits and local constants, and the orbits whose positions are good and
whose total local constant is small are the served ones.

% before: pd-gain-link
The inequality of Chapter~\ref{chapter-inert-prime-balance} needs the gains at inert primes to be
at least the class sums, up to errors small enough to be absorbed in the main term. The gain link
combines this lower bound with an upper bound on each gain, local conditions on the lattices
$N^\circ_u$ at the places where gains are counted, and conditions on the served orbits.

% before: pd-projection-property
The gains at an orbit of small primes are bounded through the lattices of the local model at a
single orbit, compared along the edges of the Bruhat-Tits building. The projection property asks for
vectors at the sources of the edges into a given lattice that are divisible along the other edges
and whose pairings with the dual lattice have nearly full rank.

% before: pd-extremal-projector
The vectors of the projection property are built from weight vectors of representations of general
linear groups. On a weight space, the extremal projector fixes the vectors killed by the raising
operators and kills those in the image of the lowering operators, so it separates the highest weight
vectors of that weight from the rest. Both gain theorems use it in this way.

% before: pd-symbols-path-generators
The symbols of the gain theorems are built from polynomials in rows of variables, one row for each
point of a segment, on which general linear groups act through each row. Such polynomials are
described by sequences of dominant weights that grow or shrink by a horizontal strip at each row, and
by products of determinants along the columns of the corresponding diagrams. Twisting each row by a
power of Frobenius gives the lifted versions that enter the symbols.

% before: pd-symbols-path-weights
The twisted degrees of a polynomial in the first rows are read off from the weights of the module it
generates under the divided powers of the polarization operators of those rows. These operators have
integer coefficients, so the module makes sense over any ring, and a lower set of weights is the kind
of condition that such a module can satisfy at its dominant weights.

% before: pd-symbols-path-count
The path generators are as many as the monomials of the same row degrees. This is the dimension count
behind Pieri's formula: each row multiplies the number of standard fillings by the number of monomials
of that row.

% before: pd-symbols-path-canonical
The independence of the path generators is proved by a highest-weight argument for the general linear
group acting on all rows at once. The weight of a path generator is read off its filling, and the
filling of largest weight for a given shape is the canonical one; acting by a matrix produces brackets of
minors, which record the filling that was replaced.

% before: pd-symbols-path-exterior
A column function is built by wedges and contractions with the rows, so the action of a matrix on it is
governed by the action on the exterior algebra. Wedging with a row is compatible with the matrix acting
on the row, and contracting is compatible with the inverse transpose; complementary minors relate the
two.

% before: pd-symbols-path-action
Acting by a matrix on all rows at once, every column function is transformed by the minors of that
matrix, or of its inverse transpose for the nonpositive columns. Expanding the product over the columns
writes the translate of a path generator through the same path with all possible fillings.

% before: pd-symbols-path-filling-weights
Each path generator has a single torus weight, determined by its filling. Among the fillings of a fixed
shape this weight is at most the shape in the dominance order, with equality only for the canonical
filling.

% before: pd-symbols-path-top-weight
Combining the last two facts, the translate of a path generator has weights at most the last vector of
its path. Its component of that top weight is the generator with the canonical filling, multiplied by the
bracket of the original filling.

% before: pd-symbols-path-canonical-independent
With the canonical filling at the end, the generators of the paths ending at a fixed vector are
independent. Reading the last row lexicographically, the leading term of such a generator is the
generator of the shorter path times a monomial that determines the previous vector of the path, which
gives an induction on the number of rows.

% before: pd-symbols-path-brackets
The brackets attached to the standard fillings of one shape are, up to a power of the determinant and a
sign, products of minors of the generic matrix along a semistandard tableau. Their independence is the
standard basis theorem for bitableaux.

% before: pd-symbols-path-independent
Over a field the path generators are linearly independent. Together with the count, this makes them a
basis over every field, and then over the integers.

% before: pd-symbols-path-basis
The path generators form a basis of the polynomials of given row degrees over every commutative ring.
This is what allows a bound on a polynomial to be read off from the paths of the generators in its
expansion.

% before: pd-symbols-path-truncate
Acting only on the first rows, a path generator is a combination of products of a generator of the
truncated path with polynomials in the later rows. The highest-weight bound for the truncated path then
bounds every weight that occurs by the vector of the path at the truncation.

% before: pd-symbols-path-lower-top
When the action reaches all the rows, the expansion of a polynomial in path generators can be read off
the dominant weights of its generic translate. A generator whose path ends outside a lower set would
produce a dominant weight outside it, by the highest-weight argument used for independence.

% before: pd-symbols-path-lower-set
Over a field, the paths that occur in the expansion of a polynomial are controlled at each row by the
module the polynomial generates under the first rows. If all the dominant weights of that module lie in
a lower set, then only paths passing through that lower set at that row occur.

% before: pd-symbols-path-translate
To compare twisted degrees with weights, one translates a polynomial by the generic invertible matrix.
The exponents of the twist are then the entries of the weights of the monomials of the translate, and
these weights are the ones the generated module sees: translation by transvections and by
diagonal matrices is expressed through the divided powers and the binomial operators.

% before: pd-symbols-path-valuation
The twisted degrees of Definition~\ref{pd-symbols-path-generators}(1) are the extreme entries of the
weights of the generated module. In the proof of Lemma~\ref{pd-symbols-path-clauses}(1) this turns
bounds on the twisted degrees into a lower set of weights, to which Lemma~\ref{pd-symbols-path-lower-set}
applies.

% before: pd-symbols-path-word-bounds
Conversely, each path generator has its twisted degrees in the first rows bounded by the last and the
first entry of its path after those rows. The bound holds over the integers, and so over every ring.

% before: pd-symbols-path-lift-twist
A lifted column function is built with Frobenius twists of a matrix at increasing positions. Undoing the
twist at one position and raising to a $p$-th power is compatible with undoing it at a later position,
modulo $p=\varpi_0^2$, because Frobenius is the $p$-th power map modulo $p$.

% before: pd-symbols-path-lift-nominal
At each row only some variables carry the uniformizer, and a costly row must use one of them. The sets of
coordinates that the steps can reach without paying are those of the plane and normal factors, and they
change from row to row in a controlled way.

% before: pd-symbols-path-lift-order
The power of the uniformizer dividing a lifted column function after the lettered substitution is counted
one costly row at a time. A step at a costly row either uses a variable carrying the uniformizer or
starts from a set of coordinates that already gained an extra factor.

% before: pd-symbols-path-lift-residue
When the substitution uses the same matrix as the lifted column function, the quotient by the counted
power of the uniformizer can be computed modulo the uniformizer. It is the Frobenius twist of the column
function built from the variables that the plane and normal factors keep.

% before: pd-symbols-path-clauses
Six facts about these polynomials are used for the symbols. The polynomials whose twisted degrees in
each initial block of rows are bounded are spanned by path generators, and a path in rank $r+K$ splits
into a path in rank $r$ and a path in rank $K$ that meet a condition at each row. Under a weighting
of the rows, the lowest part of a column function or of a path generator splits into a plane factor
and a normal factor, and the lifted column functions realize these lowest parts modulo the uniformizer;
and a twisted degree larger than the degrees of the other vector rows forces a power of the determinant
of $r$ vector rows to divide the polynomial.

% before: pd-subspace-count
The edges into a lattice correspond to subspaces of a vector space over the residue field of $W$.
Their number enters the projection property through a lower bound, and the elementary bound below is
enough.

% before: pd-edge-basis
Along an edge, the lattices at its source and target and the lattice $H(e)$ have a common basis, up
to powers of $\varpi_0$. Hence the edge pairing makes $R(e)$ and $R'(e)$ dual to each other, and an
integrality condition on pairings can be checked one edge at a time.

% before: pd-edge-transport
The group $GL_k(Q)$ acts transitively on the lattices in $Q^k$, and its action on $V$ respects the
lattices, the edges and the edge pairing. So the projection property needs to be checked at one
lattice only.

% before: pd-edge-family
The vectors of the projection property at a lattice are collected edge by edge. At a single edge they
must lie in the lattices at its source, be divisible along the other edges with the same source, and
pair with test forms at the target to divisors of a fixed power of $p$.

% before: pd-edge-assembly
Families at different edges with the same target do not interfere: a test form chosen for one edge is
divisible along the others, so it pairs with their vectors to multiples of $p^{D+1}$. A rank argument
over the discrete valuation ring $\mathcal R$ then turns the test forms into the map $u$ of the
projection property.

% before: pd-edge-count
The size of each edge family is bounded below by $\dim V$ divided by $p^{n'r(k-r)}$, so enough edges
into the standard lattice are needed to recover all of $\dim V$. Their sources can be written down as
kernels of explicit maps to $(W/pW)^r$.

% before: pd-edge-count-subspaces
The edges built in the proof of Lemma~\ref{pd-edge-count} are those whose sources reduce to graphs of
linear maps. In general the sources correspond to subspaces of the residue space.

% before: pd-edge-products
The vectors at an edge are built as products of vectors for two smaller weights. Products of the
lattices of Definition~\ref{pd-projection-property} stay in the lattices of the summed weight, so
integrality and divisibility pass to the products.

% before: pd-edge-dimension-count
At a top position the weight splits into a symmetric power of degree $d_t$ and a remaining weight for
$GL_{k-1}$. The Weyl dimension formula shows that this split loses at most a factor
$1-kc_0/a$ of the dimension at each top position.

% before: pd-edge-situation
The edge step and the numerical comparisons below share the hypotheses of
Theorem~\ref{pd-gain-theorem-two}, with a variable constant $C$ in place of $C_0$. The quantities
$t_0$ and $\omega$ measure the losses at the source and the target of an edge.

% before: pd-edge-numerics
The constants of the edge step must fit the bounds of Theorem~\ref{pd-gain-theorem-two}. The
comparisons follow from the size conditions once $C$ dominates the constants of the edge step.

% before: pd-symbols-run-system
The vectors of the edge step are built from polynomials attached to the positions of $T$, lifted
periodically to $\mathbf Z$. The lifted positions are grouped into runs of $r$ consecutive even
elements of $T$, and each run gathers the elements of $T$ since the previous run. The conditions on a
run system keep the runs short and the elements around them well separated.

% before: pd-symbols-run-system-exists
The tiling conditions of Chapter~\ref{chapter-host-fields} are what produce a run system for $T$.
The runs come from the even elements of $T$ in marked tiles, and the window conditions give the
bounds on the number and spacing of the elements of $T$ between runs.

% before: pd-symbols-module
Each run gives a module of polynomials in the variables at its elements, cut out by weight conditions
under a generic one-parameter subgroup and by degree bounds. Raising the variables to powers of $p$
moves these polynomials to coordinates indexed by $\mathbf Z/n'\mathbf Z$, and the products of the
results with small monomials are the symbols. The symbols of the right multidegree are the
candidates for the flat part of the edge family.

% before: pd-symbols-count
At each edge the edge family must reach about a $p^{-n'r(k-r)}$ share of the flat dimension. The count
of symbols gives a direct summand with a basis of symbols whose rank is this share up to a factor
close to $1$. It also bounds the normal degrees of all symbols, which later bound the levels of the
vectors.

% before: pd-symbols-space
The hypotheses of the count hold at the data of the edge step. The degrees are the gaps between the
first two entries of the weights at $T$, and the remaining integers are the ones the realization of
symbols needs.

% before: pd-symbols-edge-basis
A symbol is turned into a vector at an edge through a basis adapted to the edge. The edges with a
given source correspond to the planes spanned by $r$ rows modulo $p$ of matrices in $GL_k(W)$, and
two matrices give the same edge when their planes agree.

% before: pd-symbols-twist
The substitution of a matrix in $GL_k(W)$ twists the coordinates at each position by a Frobenius
conjugate of the matrix and scales by $\varpi_0$ the coordinates that the edge divides by $p$.
Integrality after the substitution then tests membership in the lattices of the edge. A symbol is
multiplied by a product of minors of the twisted coordinates before it is realized.

% before: pd-symbols-realize-keys
Realizability is proved for a set of polynomials built from keys: integral polynomials whose image
under the substitution of the edge agrees, to one more power of $\varpi_0$, with a polynomial of fixed
normal degree, and which are divisible to that order at every edge. Products of keys times polynomials
without normal variables span this symbol set.

% before: pd-symbols-realize-closure
The symbol set is assembled from the factors of a symbol multiplier, so it must be closed under
products. Taking the part of fixed normal degree must also stay inside it after multiplication by a
polynomial without normal variables.

% before: pd-symbols-realize-segments
The first factors of a symbol are the images of elements of the segment modules. Each of these images
lies in the symbol set.

% before: pd-symbols-realize-points
The label of a symbol has no such structure by itself, and its normal variables are treated one at a
time. A normal variable at a top position, times the determinant of that position in the symbol
multiplier, lies in the symbol set of normal degree one. The keys are bordered minors of the rows below
the position when it is even, and pairings of these rows with the dual row when it is odd.

% before: pd-symbols-realize-labels
The exponents of a label are at most $q$, so each normal variable of the label can be paired with its
own copy of the determinant of its top position in the symbol multiplier. Multiplied by these
determinants, the label becomes an element of the symbol set of its normal degree.

% before: pd-symbols-cleaning-terms
The polynomial that cleans a key product is built by an induction on terms: products of a power of
$\varpi_0$, keys and an integral polynomial without normal variables. Each term has a symbol, formed
from the second entries of its keys, and a raw polynomial, formed from the keys themselves and the
localizer of $g$. The localizer is a product of differences of powers of $r\times r$ minors, and its
part without minors on other columns is a product of powers of the determinants $\Delta_j(g)$.

% before: pd-symbols-cleaning-off-plane
The raw polynomial of a term must vanish to high order at the planes other than that of $g$, and each
of the $c'$ factors of the localizer gives one order. Modulo $\varpi_0$, a minor of the localizer at
another plane is a determinant of rows times a power of a minor of a matrix over the residue field of
$W$. The exponents in the localizer make these powers equal to $1$, so that the two products in one of
its factors agree.

% before: pd-symbols-cleaning-expansion
At the plane of $g$ itself the localizer differs from $\Delta^\alpha_g$ by terms with minors on
columns beyond $r$. After the substitution each such column brings a factor $\varpi_0$ with a normal
variable. So these terms carry at least as many factors $\varpi_0$ as normal variables, and keep all but
$e$ factors of each power of a substituted determinant, where $\varpi_0^e$ is their power of
$\varpi_0$.

% before: pd-symbols-cleaning-difference
The substituted raw polynomial of a term differs from its symbol through the tails of its keys and the
expansion of the localizer. Both contribute products with more factors $\varpi_0$ than the level of the
term, and their normal variables are absorbed by the top determinants, one key of order one for each.
So the difference is a sum of symbols of terms of higher level, up to a high power of $\varpi_0$.

% before: pd-symbols-realize-cleaning
The remaining factor of the symbol multiplier is fixed and involves no normal variables. A key product
times this factor and an integral polynomial without normal variables agrees at its plane, up to a high
power of $\varpi_0$, with the substitution of an integral polynomial that vanishes to the same order at
the other planes.

% before: pd-symbols-realize-core
The part of fixed normal degree of an element of the symbol set times the fixed factor is a sum of such
products, up to a high power of $\varpi_0$. So elements of the symbol set times this factor are
realizable at the level of their normal degree.

% before: pd-symbols-realize
The multiplier of each symbol splits into pieces, one for each level up to the normal degree of the
symbol, each integral and realizable at its level. A realizable piece agrees at its edge with an
integral polynomial up to a high power of $\varpi_0$, and that polynomial vanishes to the same order
at the other edges with the same source.

% before: pd-symbols-flat-model
The space $V^\flat$ of the flat weights is a space of polynomials of a fixed multidegree, with one
block of variables for each position. A flat model identifies the two so that integral coefficients
correspond to the lattice at the source, and integrality after the substitution of an edge to the
lattice of that edge.

% before: pd-symbols-cloud
The realized symbols cover the flat part of the weight; the remaining part needs vectors of
$V^\sharp$ and forms that detect the products. The flat model, as many weight vectors of $V^\sharp$
at the edge as the count of Lemma~\ref{pd-edge-dimension-count}, and the forms at the target all come
from the same data.

% before: pd-symbols-flat-source
Each realizable piece of a symbol multiplier gives a vector of $V^\flat$ through the flat model. The
vector lies in a bounded multiple of the lattice at the source, is close to the image of the piece at
its edge, and is small at the other edges with the same source.

% before: pd-symbols-rest-bound
The weight vectors of $V^\sharp$ lie in the lattice of the edge, while the edge family needs them in a
multiple of the lattice at the source. The entries of the weights $\lambda^\sharp$ are bounded by the
size conditions, so the two lattices differ by at most a factor $\varpi_0^{-\omega}$.

% before: pd-symbols-select
The rank in the count of symbols is computed modulo the uniformizer. Over the residue field the span of
the reduced symbols has a basis of reduced symbols, and the symbols lifting it are linearly
independent and span a saturated submodule. Over a discrete valuation ring such a submodule of a
finite free module is a direct summand.

% before: pd-symbols-normal-degree
A symbol is a product of the images of elements of the segment modules and of a label. The conditions
on the monomials of a segment module bound the normal degree of its image by a multiple of the number
of new elements of the segment, and these numbers add up to at most $|T|$. The exponents of a label
are below the twist bound $q$, which bounds its normal degree by $|T|kq$.

% before: pd-symbols-residue-coarse
The symbols are counted through simpler polynomials, built from coarser segment modules with extra
twists by powers of $p$. A product of determinants in the plane variables, the discriminant, carries
these coarse generators to symbols.

% before: pd-symbols-residue-discriminant
Multiplication by the discriminant must not lose dimension over the residue field, so the
discriminant must stay nonzero modulo the uniformizer. Each of its determinants involves rows at
distinct positions, which an evaluation detects.

% before: pd-symbols-residue-top-count
Multiplying a coarse generator by the discriminant raises its multidegree by that of the discriminant,
and the product must still have the multidegree of the symbols. So the degree of the discriminant at
each residue in $T$ must stay below the degree available there. Each residue in $T$ occurs in at most
two segments, which bounds that degree by a small fraction of $a$.

% before: pd-symbols-residue-pairs
A segment of a run system is a lettering in its own right: its points, in increasing order, are vector
letters when even and covector letters when odd. Pairs of paths for this lettering, one in the plane
directions and one in the normal directions, give products of path generators, and the conditions on
the pairs mirror the conditions on twisted degrees in the definition of a segment module.

% before: pd-symbols-residue-segment-span
The segment modules are defined by bounds on twisted degrees. They are spanned by products of path
generators, so a statement that is linear in the elements of a segment module reduces to one about
these products.

% before: pd-symbols-residue-top-degree
A nonzero coarse generator of the right multidegree uses, at each point strictly inside a segment, almost
all of the degree available at its residue. This is a lower bound for the row degrees of its segment
factors, which the lifting of the next lemma needs.

% before: pd-symbols-residue-lifts-numeric
The lift of a twisted segment generator compares the coarse twist $p^{R_j}$ with the parameters of the
symbol spaces. Three inequalities between these quantities are used, one for each condition of the
segment module that the twist must respect.

% before: pd-symbols-residue-lifts
After multiplication by the factor of the discriminant at a segment and a twist by Frobenius, a product
of path generators in a coarse segment module agrees with an element of the segment's symbol space
modulo the uniformizer. The exponents allowed by a coarse label are small enough not to disturb this.

% before: pd-symbols-residue-families
The coarse generators are counted through families of row degrees, one for each segment. A family is fixed
by the remainders of the degrees at each position modulo the coarse twists and by a choice of degrees at the
rows that link consecutive segments.

% before: pd-symbols-residue-count-digits
The families of remainder vectors split the monomials of a given degree by their digits modulo a power of $p$.
Counting the monomials by their remainders bounds a binomial coefficient by a sum over the remainders, and when the
degree is large there are many remainders.

% before: pd-symbols-residue-count-sum-facts
The summation over the families uses how the new positions of the segments tile $T$, how the twists $P_j$ and $Q_j$
fit together along the run system, and two comparisons of binomial coefficients. The comparisons hold because the
degrees $d^D(t)$ are large compared with $Y$, $k$, $N_s$ and $P_j$, and lose only factors close to $1$.

% before: pd-symbols-residue-count-inputs-products
The coarse generators are products of images of polynomials under maps that send monomials to monomials. When
distinct families of monomials have distinct images, products of linearly independent families stay linearly
independent, and families whose images lie in disjoint sets of monomials do not interact.

% before: pd-symbols-residue-count-inputs-saturated
A basis of a coarse segment module should reduce to a linearly independent family modulo the uniformizer. This
holds for a submodule that contains every polynomial of which it contains a nonzero multiple. For segment modules
it follows from the description by segment generators, which are part of a basis of the polynomial ring.

% before: pd-symbols-residue-count-inputs-reductions
With this, the coarse segment modules are free of finite rank, and their bases keep their rank modulo the
uniformizer. The ranks in the count of the coarse generators are therefore dimensions over the residue field.

% before: pd-symbols-residue-count-inputs-old
The elements of a segment that are not new are the points $f(j-1,i)$, which are new in the previous segment. The
comparison of exponents below needs their residues and the twists $P_j$ at them in terms of the new elements.

% before: pd-symbols-residue-count-inputs-monomials
A remainder vector and a family of monomials give, under the maps $\psi_j$, one monomial in the $X_{t,c}$. Its
exponent at a residue can be read as digits in base $P_j$ and $Q_j$, so the monomial determines the remainder vector,
the family of monomials and the vector $e$.

% before: pd-symbols-residue-count-inputs-index
The lower bound for the rank of a coarse segment module at one segment comes from an explicit family of its
elements. Each is a power of a determinant in the old rows times a segment generator, indexed by a path whose
degrees are fixed below $f(j,0)$ and by exponents at the rows $f(j,i)$ that stay below $Q_j$.

% before: pd-symbols-residue-count-inputs-ratio
The exponents $u$ at the rows $f(j,i)$ must leave room below $Q_j$ after the degrees fixed by the path. This room exists
because the elements below $f(j,0)$ are far from the end of the segment, so their total degree $\beta_j(b)$ is small
compared with $Q_j$.

% before: pd-symbols-residue-count-inputs-rows
The degrees of a counting path at the rows $f(j,i)$ are also controlled by $\beta_j(b)$. The path starts at $0$, its last
entries can drop only at covector letters below $f(j,0)$, and its first entries at the rows $f(j,i)$ equal $\beta_j(b)$.

% before: pd-symbols-residue-count-inputs-residues
The exponents $u$ are then counted residue class by residue class. In a box of side $Q-2\beta$ almost a proportion $1/Q$
of the integer vectors have a sum in a given class modulo $Q$.

% before: pd-symbols-residue-count-inputs-segment-clauses
The families of a segment are handled with five facts about polynomials and fillings. Twisted degrees add over
products, which places the power of the determinant $D_j$ in the segment module. The numbers of standard fillings obey
the Weyl dimension formula, the branching rule and Pieri's rule, which count the counting paths, and the products of
plane and normal path generators of a segment are linearly independent, which separates the elements of a family.

% before: pd-symbols-residue-count-inputs-segment-plane
The first $r$ entries of a counting path stay between $0$ and $\beta_j(b)$ and reach $\beta_j(b)$ at the rows $f(j,i)$.
So the canonical filling of their last vector takes only the values $\{1,\dots,r\}$ and $\emptyset$, which makes it a
standard filling and makes the plane factor of an element of a family explicit at the last $r$ letters.

% before: pd-symbols-residue-count-inputs-segment-members
The elements of a family lie in the coarse segment module because a counting path splits into a plane path and a
normal path that fit together at every letter. The power of the determinant $D_j$ raises the twisted degrees of the plane
factor by $m_0$, which is the margin the segment module asks for. The exponents at the rows $f(j-1,i)$ and $f(j,i)$ make
up the row degrees $b$, and those at the rows $f(j,i)$ stay below $Q_j$ because the path adds at most $2\beta_j(b)$ there.

% before: pd-symbols-residue-count-inputs-segment-kill
The linear independence of the elements of a family is reduced to that of simpler polynomials in two ways. Setting
some variables to $0$ keeps the part of lowest degree in them, so independence after the substitution gives
independence before it. Multiplying independent polynomials by monomials in other variables keeps them independent
as long as the pairs of a monomial and a polynomial are distinct.

% before: pd-symbols-residue-count-inputs-segment-cut
At the last $r$ letters the plane path rises by at most one entry at a time. Setting all but one plane variable of each
of these rows to $0$ turns the plane factor into a monomial times the path generator of the path stopped before these
letters. The stopped path and its filling still determine the path, so the elements of a family stay distinct after the
substitution.

% before: pd-symbols-residue-count-inputs-segment-paths
Without the conditions on the thresholds, the counting paths of a segment, with the standard fillings of their normal
parts, are counted by a product. Pieri's rule at the letters below $f(j,0)$ and the branching rule at the last $r$
letters give the product of the binomial coefficients of the prescribed degrees.

% before: pd-symbols-residue-count-inputs-segment-typical
Most of the paths in this count satisfy the conditions on the thresholds. The typical ones are those whose positive and
negative parts change little at each letter, compared with the size of its row, and whose $r$-th entry is large compared
with the degree of the oldest earlier row.

% before: pd-symbols-residue-count-inputs-segment-atypical-setting
The count of the atypical paths uses only a few properties of the segment: the sizes of the row degrees, their growth
along letters of one type, which letters are vectors, and the oldest earlier vector letter of each late letter. Stating
these as a setting separates the combinatorics of paths and the estimates of binomial coefficients from the arithmetic
of the run system.

% before: pd-symbols-residue-count-inputs-segment-atypical-rows
The arguments below compare entries of a path at different letters. A strip moves each entry by at most the size of
the strip and shifts the entries by one place, so the entries are controlled by the sums of the row degrees met so far
and by the number of vector letters passed.

% before: pd-symbols-residue-count-inputs-segment-atypical-orders
The first half of the comparison is that typical paths satisfy the thresholds. The $r$-th entry of a path is bounded by
twice the row degree of the oldest earlier vector letter, as the earlier row degrees of that type form a rapidly growing
sequence, and this absorbs the shift of a strip into the normal entries.

% before: pd-symbols-residue-count-inputs-segment-atypical-count
The atypical pairs are counted path by path. The paths themselves are few: all their entries are bounded by the total
row degree, and only the first and last $N$ entries of each vector can be nonzero.

% before: pd-symbols-residue-count-inputs-segment-atypical-tails
The standard fillings of the normal part are counted by the Weyl product of its last vector, which is compared with the
row degrees through the tails of positive and negative entries. Each tail is bounded by the corresponding tail of the
row degrees of the latest letters of one type, and a deficit in such a bound at one letter persists to the end of the
path.

% before: pd-symbols-residue-count-inputs-segment-atypical-log
The comparison is made on logarithms. The logarithm of a binomial coefficient is a concave sum of logarithms, the
logarithm of the Weyl product splits into such sums over the positive and the negative entries up to an error of order
$N^2$, and a deficit in the tails lowers the sum by an amount proportional to the deficit.

% before: pd-symbols-residue-count-inputs-segment-atypical-weight
Combining these, an atypical path has few standard fillings of its normal part. Each of the three ways of failing to be
typical produces a deficit in one tail, and the deficit lowers the logarithm of the Weyl product below that of the rank
by a fixed proportion of $K-2N$.

% before: pd-symbols-residue-count-inputs-segment-atypical-dictionary
The segment of a run system gives an atypical setting. The sizes of the row degrees come from the bounds on the
discriminant-corrected degrees and from the balance of $b$, and the growth along letters of one type from the distances
to the end of the segment.

% before: pd-symbols-residue-count-inputs-segment-atypical
The thresholds of a counting path hold along every typical path, and the atypical paths carry few standard fillings of
their normal parts. Together these show that the counting paths make up all but a proportion $a^{-N_s^2}$ of the count
without thresholds.

% before: pd-symbols-residue-count-inputs-segment
Three facts about these families give the bound. The elements are linearly independent members of the coarse
segment module, there are many counting paths, and the exponents at the rows $f(j,i)$ can be chosen in every
residue class modulo $Q_j$ in many ways.

% before: pd-symbols-residue-count-inputs
The count of the coarse generators compares the dimension of their reductions with products of the ranks of
the coarse segment modules over these families. A lower bound for each rank at one segment, and a summation
over the families, then give the bound of the statement.

% before: pd-symbols-residue-generators
Two facts about the coarse generators remain. The discriminant carries each of them to a symbol modulo
the uniformizer, and their reductions span a space of the size that the count of symbols needs.

% before: pd-symbols-residue-rank
The lower bound for the rank in the count of symbols is a statement over the residue field of
$\mathcal R$, about the span of the reductions of the symbols. Lemma~\ref{pd-symbols-select} turns it
into a direct summand over $\mathcal R$ with a basis of symbols.

% before: pd-symbols-flat-letter
At a single position the flat weight has one nonzero entry, first or last. Its dual Weyl module then
consists of the polynomials in the first column of the matrix of variables, or in the last row of its
adjugate. Twisting these coordinates by a matrix moves the lattice formed from the columns of that
matrix to the standard one.

% before: pd-symbols-flat-exists
The flat weight has at most one nonzero entry at each position. Its space consists of the polynomials in
the first column of the matrix of variables at even positions and in the last row of its adjugate at
odd positions, so substituting these, twisted by $\gamma$, for the variables $X_s$ identifies
polynomials of a fixed multidegree with it. Under this identification integrality and the lattices of
the edges become conditions on the coefficients.

% before: pd-symbols-primitive-basis
The weight vectors for the remaining part of the weight are built from a basis of the dual Weyl
modules whose leading monomials are distinct. At a top position only part of the basis is used, and
its size is the dimension for the weight with its first entry removed, which is the count of
Lemma~\ref{pd-edge-dimension-count}.

% before: pd-symbols-primitive-independent
Modulo the uniformizer, the products of the flat coordinates with the chosen basis vectors must stay
linearly independent. Their leading monomials determine both factors, because the chosen basis
vectors at a top position leave the entries read off by the flat coordinates fixed.

% before: pd-symbols-primitive-transfer
Twisting by the basis of the edge turns the lattices of the edge into lattices of integral
polynomials. Applied to a sum of products of flat vectors and weight vectors, it turns divisibility in
the lattice of the edge into divisibility of the corresponding sum of products of integral
polynomials.

% before: pd-symbols-primitive
The vectors of an edge family are products of realized symbols and weight vectors for the remaining
weight. The test forms of Lemma~\ref{pd-symbols-test} are built for combinations of these products that
are not divisible by $\varpi_0$ in the lattice of the edge, so this must hold whenever a coefficient is
a unit.

% before: pd-symbols-test-operators
The forms at the target are first built on the flat part of the weight, for polynomials obtained from
the realized symbols by a bounded number of first-order operators at the top positions. These
operators move a combination of symbols to one whose image in the flat model is not divisible by the
uniformizer.

% before: pd-symbols-test-size
The map of the next lemma loses a power of $p$ that comes from comparing dual lattices of the dual Weyl
modules. Its exponent is a sum over the orbits of a logarithm of the spread of the weight, so it is
small compared with the degree of the weight.

% before: pd-symbols-test-onto
The weight at the target is the sum of the flat weight and the weight of the remaining factors. An element
of the lattice of an edge for this sum is spanned by products of elements of the lattices for the two
summands, which lets a form on the whole space be read through products.

% before: pd-symbols-test-fock
The identifications of the next lemmas pair a polynomial with another through their coefficients, weighted
by the factorials of the exponents. This pairing is multiplicative over disjoint sets of variables, so on a
product of dual Weyl modules it is the product of the pairings at the positions.

% before: pd-symbols-test-fock-position
At one position the Fock pairing, rescaled, compares the lattice of a dual Weyl module with the one of the
dual lattice. Integrality holds after multiplying by the power $p^{N(\mu)}$, whose exponent is the size bound
of Lemma~\ref{pd-symbols-test-size} at that position; the comparison for any lattice follows from the
standard one by a change of basis.

% before: pd-symbols-test-fock-dual
Tensoring the comparisons at the positions identifies the whole space with its dual through the Fock
pairing. The same bases describe the lattices of a lattice and of an edge, so the identification moves the
dual lattices into the lattices of the dual lattice and the dual edge, up to $p^{\kappa_1(\lambda_0)}$.

% before: pd-symbols-fock-words-operators-symbol
The bound on products of the operators $Y_a\,\partial/\partial Y_b$ is proved on symbols. An endomorphism of the
homogeneous polynomials of degree $d$ is recorded by a polynomial in two sets of $k$ variables, of degree $d$ in each,
and the commutators with the operators become derivations of this polynomial ring.

% before: pd-symbols-fock-words-operators-kappa
The symbol map is injective, and it turns composition with $Y_a\,\partial/\partial Y_b$ on the left and on the right into
explicit first-order operators. The identity goes to a multiple of the $d$-th power of $\sum_iP_iQ_i$.

% before: pd-symbols-fock-words-operators-tau
The image of the space of endomorphisms will be divided by powers of $\tau=\sum_iP_iQ_i$. The derivations $R_{ab}$
commute with multiplication by $\tau$, eigenvectors of $R_{aa}$ have one weight at $a$, and a bihomogeneous polynomial
that vanishes on the quadric $\tau=0$ is divisible by $\tau$.

% before: pd-symbols-fock-words-operators-exponential
For $a\neq b$ the derivation $R_{ab}$ is locally nilpotent. Its exponential is a linear change of coordinates that
preserves the quadric, which moves points of the quadric to simpler ones.

% before: pd-symbols-fock-words-operators-quadric
A space of symbols that is stable under the $R_{ab}$ and has small weights vanishes on the quadric. At a point with one
nonzero coordinate among the $P_i$ the weights are too large for a monomial to survive, and the exponentials of the
$R_{ab}$ reach every other point of the quadric from such a point.

% before: pd-symbols-fock-words-operators-words
Conversely, every multiple of $\tau^{d-c}$ by a bihomogeneous polynomial of degree $c$ in each set of variables is the
symbol of a combination of products of at most $c$ operators. Multiplying by $P_aQ_b$ costs one operator.

% before: pd-symbols-fock-words-operators
The operators $Y_a\,\partial/\partial Y_b$ span the action of $\mathfrak{gl}_k$ on homogeneous polynomials. A space of
endomorphisms that is stable under commutators with them and is spanned by vectors of weights at most $c$ consists
of combinations of products of at most $c$ of these operators.

% before: pd-symbols-fock-words-position
At one position, multiplication by a fixed element of a dual Weyl module is moved across the Fock pairing as an
endomorphism of the homogeneous polynomials. These endomorphisms are stable under commutators and have weights
bounded by the first entry of the weight, so they are combinations of products of at most $c$ operators.

% before: pd-symbols-test-fock-words
Multiplication by a fixed vector of the remaining factors must be moved across the Fock pairing onto the
flat side. On the flat variables it becomes a combination of first-order operators, a bounded number at each
top position, because the remaining weights at the top positions are bounded by $c_0$.

% before: pd-symbols-test-duality
To move forms between the flat space and the whole space, each space is identified with its dual, up to
the power of $p$ of the size bound, on all lattices at once. The product with a flat element is then
matched by an operator word in the flat variables, which keeps the forms in the operator closure.

% before: pd-symbols-test-maps
The forms at the target are obtained by composing with linear maps from the whole space to the flat space.
These maps must keep the lattices of the edges up to a bounded loss, turn products with a flat element into
images of operator words, and detect divisibility by the uniformizer.

% before: pd-symbols-test-maps-top
On the top positions the maps come from the two identifications of Lemma~\ref{pd-symbols-test-duality}: a
form on the flat space is moved to the flat space, multiplied by a remaining vector and moved back. The loss
$r|T|c_0$ comes from the remaining weight at the top positions.

% before: pd-symbols-test-maps-peel
Outside the top positions the flat weight vanishes, and the remaining weight is the whole weight there. Those
positions are removed by pairing their factors with forms of the remaining weight, which costs at most the
spread of that weight, $2rn'c_{pp}$ in the situation.

% before: pd-symbols-test-maps-cloud
A family of such maps produces the polynomial, the map on forms and the scalar of
Lemma~\ref{pd-symbols-test-cloud}. A map that does not divide the image of the combination by too high a power
of $\varpi_0$ gives the polynomial, and its dual gives the map on forms.

% before: pd-symbols-test-cloud
A combination of products that is not divisible by the uniformizer at the edge is traced back to a
polynomial in the operator closure of the realized symbols. A linear map carries forms on the flat
space to forms on the whole space, losing at most a factor $\varpi_0^{-\omega}$ at every edge and
multiplying the pairing by a bounded power of $p$.

% before: pd-symbols-test-one-sided
On the flat space the Fock pairing of an element of the operator closure is tested against products of
segment polynomials of the symbol spaces, with lowered weight bounds and degrees near those of the symbols.
These products are built only from segment modules and labels, so their realizations can be controlled.

% before: pd-symbols-test-chains
The test elements are found by a projection argument in a second polynomial ring, indexed by the letters
and links of chains of positions. The argument uses valuation bounds along each chain, a Fock pairing twisted
by units, and a pair of polynomials with prescribed valuations, which the projection produces.

% before: pd-symbols-one-sided-real-fixed
A one-sided element carries no factor of the fixed symbol multiplier, so the factor has to come out of its segment
parts. Each segment part will be divisible by a product of determinants of rows of the segment, and modulo
$p^{c'}$ the images of these determinants multiply to the fixed factor times a polynomial without normal variables.

% before: pd-symbols-division-rows
The determinants that divide a segment part are formed from the rows $a_l(x)$ below each new position $x$ and
from the rows $f(j-1,l)$ that are not new. These rows lie in the segment, and every other even row below $x$ lies
below $\operatorname{old}_j(x)$, where the degree bound on a one-sided element applies.

% before: pd-symbols-division-numeric
Each division by a determinant uses up one unit of the valuations of a segment part. The total number of
determinants is at most the margin $m_*$, and $m_*$ is at most the slack $\lfloor B_j(x)/(Yk^2)\rfloor$ built into
$\theta^1_j$; both are numerical consequences of the situation.

% before: pd-symbols-division-domain
Divisibility by a determinant over a field gives divisibility over a domain, as the determinant has coefficients
$\pm1$. Applied to several determinants in turn, it removes each of them and lowers the valuations by at most the
total exponent.

% before: pd-symbols-one-sided-real-division
The margin $m_*$ of a one-sided element is large enough to divide each product of its segment module by these
determinants. What is left is a generator of the segment module with margin $0$, whose image is a symbol.

% before: pd-symbols-test-one-sided-real
To turn a test element into a form on the flat space it must be realizable at the edge, with the margin
$c'$. For one-sided elements this holds at every $g$.

% before: pd-symbols-flat-real-multinomial
The forms on the flat space are Fock pairings divided by $\prod_sd(s)!$, so their integrality needs a bound
for multinomial coefficients. At each position the power of $p$ in such a coefficient is at most $k$ times
$\lfloor\log_pd(s)\rfloor$, which gives the factor $p^{\kappa_F}$.

% before: pd-symbols-flat-real-forms
A polynomial of multidegree $d$ gives a form on the flat space by its Fock pairing with the inverse of the flat
model at $1$. A linear change of variables is adjoint, for the Fock pairing, to the transposed change, and the
lattices of a $W$-lattice and of an edge are the images of the integral polynomials under such changes. So the
form of an element of the lattice of the dual is integral on the lattice itself.

% before: pd-symbols-test-flat-real
A realizable polynomial of the right multidegree gives a form on the flat space with the lattice properties
required at the target. Its pairing with an element of the flat model is, up to $\varpi_0^{c'}$, a Fock
pairing on the side of the twisted substitution.

% before: pd-symbols-test-flat-chains-tied
The test elements for a polynomial $G$ of the operator closure come from signed chain data whose letters of level $j$
are the elements of $\Lambda'_j$ and whose links are the $f(j,i)$. Their thresholds are formed from the degrees $B_j$,
the integers $\lfloor B_j/(Yk^2)\rfloor$ and $m_*$, and weights of the normal variables. A polynomial on the positions of
such a datum is read off from $G$ by keeping the monomials with prescribed residues, dividing their exponents and
projecting the normal part onto the weights.

% before: pd-symbols-test-flat-chains-membership-loss
An operator of the closure lowers the twisted degrees of the plane coefficients at the positions of the residue it
acts on. This is recorded by subtracting a loss from the thresholds of the chain data. Each operator costs at most $2$
at an index, and only at the indices whose place is at least some element of the level with its residue.

% before: pd-symbols-test-flat-chains-membership-one-position
An operator $X_{s,a}\,\partial/\partial X_{s,b}$ acts on the variables of one residue $s$. For a chain datum tied to the
run system, exactly one position has that residue, and the cost $c_s$ is $2$ at every index whose chain or link set
contains it.

% before: pd-symbols-test-flat-chains-membership-one-transport
A derivative or a multiplication by $X_{s,c}$ changes the exponent at one place, and so its residue modulo $\lambda(y_0)$
or its quotient. On the level images this becomes a combination of the identity, an Euler operator, a derivative and a
multiplication by the variable of the level image, with the residue or the block shifted by one.

% before: pd-symbols-test-flat-chains-membership-one-twist
The twisted degrees $\operatorname{val}_S$ and $\lambda_S$ are formed by a substitution in the rows of $S$. A derivative
or a multiplication by a variable of one row moves them by at most one, and only when the row lies in $S$.

% before: pd-symbols-test-flat-chains-membership-one-plane
On the plane side, the operators of the previous lemmas lower the twisted degrees of an admissible polynomial at the
chains and link sets that contain the position by at most one. Raising those thresholds by one beforehand keeps the
result admissible.

% before: pd-symbols-test-flat-chains-membership-one-normal
On the normal side, the operators move the highest weights of the projected components by at most one at the sets
$N_j(u)$ that contain the position. Components whose weights are too low are killed by the projector, and the others are
controlled by the hypothesis for nearby weights.

% before: pd-symbols-test-flat-chains-membership-one
The membership condition with a loss passes through one operator of the closure at the price of its cost. The
derivative and the multiplication by a variable each lower the thresholds at the chains and links that the position of
its residue touches by at most $1$.

% before: pd-symbols-test-flat-chains-membership-bound
Each residue occurs in at most $c_0$ operators of a word, and the elements of a level are at most $N_1$ in number. The
size conditions of the situation make $16N_1c_0$ smaller than the integers $\lfloor B_j/(Yk^2)\rfloor$, so the loss of a
word stays below an eighth of them.

% before: pd-symbols-test-flat-chains-membership-words
With the cost of one operator and the bound for a word, the margins $(0,\frac18)$ become $(\frac18,\frac38)$ after a
word of operators. The loss of the word is absorbed into the larger margins.

% before: pd-symbols-test-flat-chains-membership-monomials-pairs-even
The polynomial $S$ is a product of determinants of matrices of powers of plane variables at even places. Its monomials
therefore add nothing to the normal variables or to the odd residues.

% before: pd-symbols-test-flat-chains-membership-monomials-pairs-image
For a product of pairs over the levels, a label and such a monomial, the level image separates into a plane part and a
normal part. The factors $\phi_j$ contribute multiples of $\lambda(y)$ to the exponents at a position, so only the label and
the monomial decide the residues.

% before: pd-symbols-test-flat-chains-membership-monomials-pairs-normal
A joint isotypic projector that does not kill the normal part forces its weight to respect the twisted degrees of the normal
factors. The first entries of the weight are bounded by the largest twisted degrees of the $G_j$ below and at each place.

% before: pd-symbols-test-flat-chains-membership-monomials-pairs-twisted
The bounds on the weight, together with the inequalities that define the segment modules, bound the twisted degrees of the
plane factors from below. These lower bounds are the numbers $\theta^*_j(u)$ of the ties.

% before: pd-symbols-test-flat-chains-membership-monomials-pairs-plane
Lower bounds for the twisted degrees of the plane factors below each place give the conditions at the chains of the plane level
image. At the links they hold up to a loss from the degrees of the links of the previous level.

% before: pd-symbols-test-flat-chains-membership-monomials-pairs-numbers
The loss at the links is small compared with the numbers $s_j(u)$. This uses the size conditions of the data and the
identities between the residues, the block degrees and the multidegree.

% before: pd-symbols-test-flat-chains-membership-monomials-pairs
The symbol $P_b$ is a product over the levels of elements of the symbol spaces, which are spanned by products $A_jG_j$
of a plane and a normal factor with bounded twisted degrees. For such products, with a label and a monomial of $S$, the
membership condition is checked at the multidegrees that the polynomial can have.

% before: pd-symbols-test-flat-chains-membership-monomials
The symbol $P_bS$ is a combination of the polynomials $P_bX^\beta$ for the monomials $X^\beta$ of $S$. For these the
level image factors over the levels, and the thresholds at the chains and links are met with the margins
$(0,\frac18)$.

% before: pd-symbols-test-flat-chains-membership
The polynomials of the operator closure satisfy the membership condition. The condition is linear in the polynomial,
so it suffices to know it for the symbols $P_bS$ with smaller margins, and to bound the loss of margin under the
operator words that generate the closure.

% before: pd-symbols-test-flat-chains-datum-assembly
A polynomial on the positions of a chain datum is turned back into a polynomial in the variables $X_{s,c}$ by raising
each variable to the power $\lambda(y)$ at the residue of its place and multiplying by a monomial of residues. The data
that are used for this are those whose degrees split the multidegree $d_T$ in blocks, and whose normal part has bounded
twisted degrees.

% before: pd-symbols-test-flat-chains-datum-pick-weights
Splitting an exponent into its residue modulo $\lambda(y)$ and its quotient splits its factorial into a factorial of the
residue, a factorial of the quotient, a power of $\lambda(y)!$ and a ratio of factorials. These ratios, and a cost of the
normal projectors, enter the comparison of Fock pairings.

% before: pd-symbols-test-flat-chains-datum-pick-inputs
The choice of the datum attached to $G$ rests on five facts. They express the Fock pairing with an assembly through the
level images, make the factorial ratios units, provide a projector and a normal test, and bound the valuations that
occur by a multiple of $|T|k^2\log_pa$.

% before: pd-symbols-test-flat-chains-datum-pick
The datum attached to $G$ is chosen so that the Fock pairing of $G$ with assemblies is, up to a scalar of bounded
valuation, the pairing $B_{\mathcal D}$ with an admissible polynomial. The error is small compared with the integer
$c'$.

% before: pd-symbols-test-flat-chains-datum-membership-levels
Each variable of the model positions and of the normal ring sits at a level and at a place of the run system. The
assembly substitutes for it a power of $X_{\bar x,c}$ fixed by the distance of the place to the top of its level. Sorting
the monomials by their degrees at the pairs of a level and a place gives the components that become one-sided
elements.

% before: pd-symbols-test-flat-chains-datum-membership-parts
Four facts turn the assembly into one-sided elements. It equals the product of the level maps and the label of the
residues, the label lies on $T$, each component has the right multidegree and one normal degree, and each component is
a sum of products of elements of the segment modules with the degree bounds.

% before: pd-symbols-test-flat-chains-datum-membership
The bounds that a projection pair satisfies at the chains and at the links are those that one-sided elements need.
So the assembly of the polynomial of a projection pair is a sum of one-sided elements.

% before: pd-symbols-test-flat-chains-datum-numbers-places
The positions of a chain datum tied to a run system have places in consecutive segments of one period. Their residues
are therefore distinct elements of $T$, which bounds the number of positions by $|T|$. The places that index the chains
are new elements of their level.

% before: pd-symbols-test-flat-chains-datum-numbers-letters
The letters of a level sit below $f(j,0)$ with a gap of more than $L_2$, so the powers $\lambda(y)$ at the letters are
large. When $\lambda(y)d(y)$ is bounded by $2a$, the degrees of the letters form a geometric series, and their sum is
small compared with $a$.

% before: pd-symbols-test-flat-chains-datum-numbers-sizes
The numbers $D_m$, $D_l$ and $\kappa_{\mathcal D}$ of a chain datum are sums of logarithms of its degrees over its
chains, links and positions. When the degrees are at most $2a$ and the positions few, each logarithm is at most a
multiple of $\log_pa$.

% before: pd-symbols-test-flat-chains-datum-numbers-slack
Two inequalities between the integers $\lfloor B/(Yk^2)\rfloor$ and the other sizes follow from the size conditions of
the situation. They bound these integers below by a combination of $c'$ and $c_0$ and above by $B/100$, and compare
them with the sum of the degrees of the letters of a level.

% before: pd-symbols-test-flat-chains-datum-numbers
The integer $c'$ is large compared with $|T|k^2\log_pa$ by the size conditions of the situation. The numbers $D_m$,
$D_l$, $\kappa_{\mathcal D}$ and $\kappa_F$ of an adapted datum are bounded by multiples of $|T|k^2\log_pa$, which gives
the hypotheses of the projection lemma.

% before: pd-symbols-test-flat-chains-datum
The datum attached to $G$ uses only the integrality of $G$ and the membership condition. Under these hypotheses there are a signed chain datum and an admissible polynomial of Fock valuation $0$ whose
projection pairs yield one-sided elements with a bounded pairing.

% before: pd-symbols-test-flat-chains
For a polynomial of the operator closure that is integral but not divisible by $\varpi_0$, the test elements
come from a signed chain datum attached to it. Every projection pair for that datum yields one-sided
elements whose Fock pairing with the polynomial divides a bounded power of $p$.

% before: pd-symbols-test-chains-projection-clauses-mahler
One of the properties used below expands ratios of factorials, as functions of an integer argument, in binomial
coefficients of that argument. The coefficients of such an expansion are the iterated differences at $0$. Their
$p$-adic valuations control how far a truncated expansion is from the function.

% before: pd-symbols-test-chains-projection-clauses-mahler-rules
Lower bounds on the valuations of these coefficients, growing linearly with their index, pass to the values of the
function, to products and to inverses. Products are handled by the Leibniz rule for the difference operator, and
inverses by a truncated geometric series.

% before: pd-symbols-test-chains-projection-clauses-mahler-factorials
The factorial ratios $c$ and $d$ of the next lemma satisfy such bounds. For $\rho<P$ the product $d$ is a binomial
coefficient whose expansion has coefficients divisible by growing powers of $p$, and $c$ is a running product of
values of $d$.

% before: pd-symbols-test-chains-projection-clauses
The construction of a projection pair uses five properties of polynomial rings over $\mathcal R$ and its fraction
field. They concern the integrality of products of isotypic projections, the description of twisted degrees by weights
and by isotypic projections, the basis of path generators, $p$-adic expansions of factorials and binomial coefficients,
and the effect of divided derivatives on twisted degrees.

% before: pd-symbols-test-chains-projection-properties
The construction itself uses the clauses only through five lifting properties of the fraction field with its
valuation. They bound the twisted degrees of variables and of divided derivatives, give a basis of products of columns
adapted to twisted degrees, give Fock projectors onto the polynomials of bounded twisted degree with a bounded loss of
integrality, and expand ratios of factorials and divided derivatives of powers in binomial coefficients.

% before: pd-symbols-test-chains-projection-ring-columns
The basis of the spanning property is made of the path generators, and its factors are their column functions. Each
column involves one variable from the row of each letter at which its height changes, and its twisted degree at an
initial set of letters is read off from its heights there. Summed over the columns, these lower bounds give the last
entry of the path.

% before: pd-symbols-test-chains-projection-ring-projectors
A Fock projector onto the polynomials of bounded twisted degree is obtained by keeping the isotypic components of
large enough last weight. The isotypic description of twisted degrees shows that it has the required image, and the
self-adjointness of the isotypic projectors makes it orthogonal. The clauses on chains of letters give its
commutations and its loss of integrality.

% before: pd-symbols-test-chains-projection-ring
The lifting properties hold over the rings of the projection lemma. They are derived from the clauses of
Lemmas~\ref{pd-symbols-test-chains-projection-clauses} and~\ref{pd-symbols-path-clauses} on polynomial rings over such
rings, through the columns of the path generators and the Fock projectors built from isotypic projectors.

% before: pd-symbols-test-chains-projection-fock-product
Thresholds on twisted degrees for several sets of letters can be imposed at once when the sets are nested or
disjoint. Fock projectors for such sets commute, so their product meets every threshold, keeps the Fock pairing with
the polynomials that already meet them, and loses a bounded amount of integrality.

% before: pd-symbols-test-chains-projection-diagonal
Multiplying each coefficient by a combination of binomial coefficients in its exponents lowers twisted degrees by at
most the total degree of those binomial coefficients. Such multipliers approximate the inverses of the units $\gamma$
of a signed chain datum.

% before: pd-symbols-test-chains-projection-lift
The first step of the construction produces an integral polynomial that meets the thresholds of the links up to $6c$
and whose pairing with $F$ has bounded valuation. It projects a monomial of $F$ of Fock valuation $0$ onto the link
conditions and corrects for the units $\gamma$.

% before: pd-symbols-test-chains-projection-good
Polynomials on the model positions are compared with polynomials on the positions through $\varphi$. The good ones
have the row degrees that $\varphi$ turns into the degrees $d(y)$, label exponents below $P(j,m)$ and tops of bounded
degree.

% before: pd-symbols-test-chains-projection-model-level-val
The moves on the model positions rename variables, multiply by monomials and extract coefficients. The twisted degrees
$\operatorname{val}_S$ behave simply under each of these, as the substitution that defines them acts on the rows of $S$
alone.

% before: pd-symbols-test-chains-projection-model-level
On the model positions, a polynomial whose link degrees all sit on the labels is moved, one level at a time, to one
whose link degrees are shared between labels and tops. The thresholds for the letters and the labels become thresholds
for the letters and the tops, and the other levels are not affected.

% before: pd-symbols-test-chains-projection-model
The integral polynomial of the first step is the image under $\varphi$ of a good integral polynomial on the model
positions. The thresholds of the links pass to the corresponding sets of model positions, level by level, and the
exponents of the labels are then reduced modulo $P(j,m)$ by moving multiples of $P(j,m)$ to the lifts.

% before: pd-symbols-test-chains-projection-omega-blocks
The pairing with $F$ sees a monomial on the model positions through $\varphi$, so through its collapsed exponent and the
units of the factorials of its label and lift exponents. A projection may change a monomial only within its block: its
row degrees and its label and top exponents stay fixed.

% before: pd-symbols-test-chains-projection-omega-cost
Twisting by units and their inverses is done with multipliers that are combinations of binomial coefficients in the
exponents. When the label exponents are fixed residues, the exponent of a link is affine in that of its lift, and such
multipliers still lower twisted degrees by at most their degree.

% before: pd-symbols-test-chains-projection-omega-units
The digit units of a monomial depend on its label exponents. For each choice of residues the units and their inverses
are approximated to order $c$ by multipliers of degree $12c$, and on sets without labels the label exponents can be
treated as fixed.

% before: pd-symbols-test-chains-projection-omega-projection
The chain conditions are imposed by the Fock projectors of property (P), row degree by row degree. On model positions
they keep labels and tops in place, as these lie in no chain.

% before: pd-symbols-test-chains-projection-omega-expansion
The change of the pairing is controlled by the digit parts of the divided derivatives of $F$ by the top exponents of a
monomial. Property (E) expands such a derivative in digits, and the hypothesis (H) absorbs the losses at the lifts.

% before: pd-symbols-test-chains-projection-omega
The chain conditions are imposed on the model positions in a last step. It replaces a good polynomial by one that
meets the chain thresholds up to $12c$, loses at most $24c$ at the link thresholds, and changes the pairing with $F$ by
an element of large valuation.

% before: pd-symbols-test-chains-projection
Lemma~\ref{pd-symbols-test-flat-chains} needs a projection pair for its datum. Such pairs exist under the hypothesis
(H) and the bound on $c$ that it provides, over every complete discrete valuation ring of characteristic
$0$ whose residue field contains the residue field of a Hida setting.

% before: pd-symbols-test-flat-xi
Combining the two lemmas gives test elements for every polynomial of the operator closure that is not
divisible by $\varpi_0$ in the flat model. Their normal degrees can be bounded, because the polynomials of
the closure have bounded normal degree and the Fock pairing only pairs equal monomials.

% before: pd-symbols-test-flat
On the flat space the forms at the target are built directly for the polynomials of the operator
closure. They are integral up to a bounded power of the uniformizer at the target, small at the other
edges into it, and their pairing with the given vector divides a bounded power of $p$.

% before: pd-symbols-test
Once a combination of the products is not divisible by $\varpi_0$ in the lattice of the edge, a form at
the target detects it. The form is integral at all edges into the target, small at the other ones,
and its pairing with the combination divides a power of $p$ bounded in terms of $|T|$, $k$ and
$\log_pa$.

% before: pd-edge-step
At each edge into the standard lattice the vectors are products of realized symbols for the flat part
of the weight and of weight vectors for the remaining part. Each edge carries about a
$p^{-n'r(k-r)}$ share of the flat dimension, which the $p^{n'r(k-r)}$ edges of
Lemma~\ref{pd-edge-count} make up, and the remaining part loses the factor of
Lemma~\ref{pd-edge-dimension-count}.

% before: pd-gain-theorem-two
For primes below $p_0$ the projection property is proved with Hecke edges of index $p^{n'r}$, for
orbits whose top positions are well spread. The defect $\kappa$ of the rank is of size $|T|r^2k/Y$,
and $t$ collects the losses from the top positions.

% before: pd-kappa-split-rank
For edges of index $p^{n'}$ the projection property is reached one edge at a time. At each edge the
vectors are taken in two small modules, the cores, whose elements are highly divisible along the
other edges with the same source or target, and the pairing matrix of a family is asked to be
invertible up to a bounded power of $p$. Flat weights, with entries only in the first and last
places at the top positions, are the weights for which these ranks are estimated.

% before: pd-kappa-edge-count
The ranks at single edges add up to a rank for the whole projection property only if there are
enough edges at each lattice. For edges of index $p^{n'}$ the edges at a lattice correspond to the
lines or the hyperplanes of a vector space of dimension $k$ over the residue field.

% before: pd-kappa-rank-monotone
The edge rank condition becomes weaker when the level $t$ or the depth grows and when the required
size shrinks. Because sizes are integers, the least defect for which it holds is attained, so the
defects defined below satisfy the condition themselves.

% before: pd-kappa-rank-criterion
When the margin $c$ is large compared with the depth, the families at different edges into one
lattice hardly interact. Their pairing matrices then assemble into one matrix congruent to the
identity modulo $\varpi_0$, which gives the projection property with the defect of the edge rank
condition.

% before: pd-kappa-rank-upper
The same assembly at the edges out of one lattice bounds the size of a family from above. The
families at distinct edges are then linearly independent, so their sizes cannot exceed the
dimension divided by the number of edges.

% before: pd-kappa-edge-move
The edges of index $p^{n'}$ are permuted transitively by $GL_k(Q)$, and the action of this group
carries the lattices that define the cores into one another. A family at one edge therefore gives a
family at every edge, so the ranks need to be produced at a single edge.

% before: pd-kappa-setting-change
The function $\kappa$ will be defined from one unramified quadratic Hida setting, while the theorem
asserts the projection property in all of them. Any two such settings are isomorphic, by the
uniqueness of Witt vectors and of unramified extensions.

% before: pd-kappa-weight-tensor
The weights needed later have small entries at the positions outside $T$ and the top part at the
positions in $T$. The space of such a weight is the tensor product of the spaces of these two parts,
compatibly with the lattices and the pairings, so the two kinds of positions can be treated
separately.

% before: pd-kappa-weight-order
For a weight whose entries are all at most $c$ in absolute value, the lattices $R(e)$ and $R'(e)$
differ from the lattices at the source and at the target of $e$ by at most $2n'c$ powers of
$\varpi_0$. This is the cost in level and margin of adding the positions outside $T$.

% before: pd-kappa-cloud-dimensions
At a top position, a weight with large outer entries and small middle entries has dimension close
to the product of the dimension of its middle part and that of the flat weight with the same outer
entries. The ratio is bounded in terms of the ratio of the middle entries to the outer ones.

% before: pd-kappa-cloud-cartan
The maps between flat weights and weights with small middle entries are built one position at a
time, on the lattices of dual Weyl modules attached to a lattice in $K^k$. Two properties of these
lattices are used: their products span the lattice of the sum of the weights, and a product is
divisible by the uniformizer only if one of its factors is.

% before: pd-kappa-cloud-primitive
Multiplication by an element that is not divisible by the uniformizer therefore detects membership in
the lattices. A retraction of lattices then gives a left inverse of this multiplication that is
integral on them.

% before: pd-kappa-cloud-transpose
The contraction at one top position is built from a pairing of $GL_{m-1}$ between the middle part of a
weight and the middle part of its dual. Its first ingredient is a map from the dual Weyl module of the
dual middle part to that of the middle part shifted to end in $0$. It is onto on the modules over
$\mathcal R$, and it is equivariant for the action twisted by the transpose.

% before: pd-kappa-gram-fock
The invariant form of the second ingredient is the Fock pairing, normalized by the value at the
leading-minor polynomial $h_\nu$. Moving a matrix from one side of the pairing to the other transposes
it, and the same holds for the divided powers of the row operations. On a dual Weyl module the pairing
with $h_\nu$ is the value at the identity matrix, and the coefficients at the lower triangular
monomials are such values; these coefficients decide membership in the module over a subring.

% before: pd-kappa-defects-gl-modules
The power of $p$ that carries the dual Weyl module over $\mathbf Z_p$ into the Weyl module comes from
a statement about representations of $\mathfrak{gl}_n$ over $\mathbf Q_p$ with a lattice. A
representation here is given by operators with the commutation relations of the matrix units and a
decomposition into weight spaces. The lattices are stable under the divided powers of the operators off
the diagonal.

% before: pd-kappa-defects-gl-highest
The bound uses standard facts of highest weight theory in characteristic $0$. The submodule generated
by a highest weight vector has its weights below that weight and meets that weight space in a line. A
finite-dimensional representation is generated by its highest weight vectors, and two highest weight
vectors of the same weight generate isomorphic submodules.

% before: pd-kappa-defects-projector
The bound comes from an explicit polynomial in the operators that projects onto the vectors killed by
all of them. It is built one index at a time: the factor for an index $z$ is a sum of iterated
conjugates $E_{zi}(\cdot)E_{iz}$ over the indices $i$ still to be treated, with coefficients for which
a sum telescopes. The same construction applies to spaces of linear maps, on which the operators act by
commutators.

% before: pd-kappa-defects-projector-invariant
On a space where the integer eigenvalues of the diagonal operators are at most $l$ in absolute value,
the polynomial with cutoff $l$ carries a vector killed by the diagonal operators to one killed by all
the operators. It changes the vector only by an element of the images of the operators off the
diagonal, and it commutes with every linear map that intertwines two such families of operators.

% before: pd-kappa-defects-projector-integral
The coefficients of the polynomial have denominators, but on a lattice stable under the divided powers
they cost only a bounded power of $p$. The binomial coefficients that appear have $p$-adic valuation at
most $\lfloor\log_p(l+k-1)\rfloor$, and there is one such factor for each of the $\binom k2$ pairs of
indices.

% before: pd-kappa-defects-schur
On the endomorphisms of the submodule generated by a highest weight vector, the polynomial is the
projection onto the scalars, given by the trace. The bound on the weights of the representation bounds
the eigenvalues of the diagonal operators on these endomorphisms, so the cutoff $l$ suffices.

% before: pd-kappa-defects-weyl-integral
Taking the component of highest weight and then applying the polynomial gives the inverse of the
dimension of an irreducible submodule times the identity, on the part of the representation generated
by the highest weight vectors of that weight. The component of a lattice vector lies in the lattice
and the dimension is an integer, so this bounds the power of $p$ that carries the lattice into the
divided-power hull of its highest weight vectors.

% before: pd-kappa-defects-sandwich-core
The integrality of the normalized pairing needs a power of $p$ that carries the dual Weyl module
into the Weyl module. Over $\mathbf Z_p$ the exponent depends only on how far apart the row degrees of
the monomials in the dual Weyl module can be, through $\lfloor\log_p(l+r-1)\rfloor$.

% before: pd-kappa-defects-sandwich
The row degrees of the monomials of a bideterminant lie between the last and the first entry of its
shape, so the bound over $\mathbf Z_p$ applies to every dual Weyl module. It passes to $\mathbf Z$ up
to a factor prime to $p$, and from there to every local domain in which $p$ is not a unit, such as
the ring of a Hida setting.

% before: pd-kappa-gram-form
With these facts the normalized Fock pairing on a dual Weyl module over $\mathcal K$ is nondegenerate,
and a vector that pairs integrally with the module over $\mathcal R$ lies in it. The Weyl module over
$\mathcal R$ pairs integrally with the dual Weyl module, and it contains $p^w$ times the dual Weyl
module, so $p^w$ times the pairing is integral on the module over $\mathcal R$.

% before: pd-kappa-gram-divisors
A form with these three properties relative to a lattice over a discrete valuation ring can be
diagonalized by elementary divisors. The diagonal entries are powers of $\pi$ with exponents at most
$N$, because the dual lattice lies between $\pi^N$ times the lattice and the lattice.

% before: pd-kappa-cloud-gram
The second ingredient is an invariant bilinear form on a dual Weyl module of $GL_r$. It is integral on
the module over $\mathcal R$, and on suitable integral elements it is diagonal with entries powers of
$p$ whose exponents are of the size of $r^2$ times the logarithm of the weight.

% before: pd-kappa-cloud-lift
Elements of the dual Weyl modules of $GL_r$ have to be lifted to $GL_{r+2}$, where $GL_r$ sits in the
middle block. A minor on the first columns lifts to the minor with one more row and one more column,
and a determinant factor splits off when the weight is shifted by a constant.

% before: pd-kappa-cloud-pairing
Dividing by a power of the determinant and composing the form with the identification of the dual
part gives the pairing. It transforms under $GL_r$ by a power of the determinant, it is integral,
and on integral elements it is diagonal with entries powers of $p$.

% before: pd-kappa-cloud-substitution
The contraction substitutes the middle block of a matrix of variables and applies the pairing to the
coefficients. Coefficients of dual Weyl modules stay in dual Weyl modules of the middle block, the
flat dual Weyl modules only pick up a power of the determinant, and a flat polynomial with the extreme
degrees in the first and the last rows is a multiple of the leading-minor polynomial $h_\nu$.

% before: pd-kappa-cloud-contraction
At one top position, a weight with outer entries $P+c$ and $-Q-c$ and small middle entries is the sum
of the flat weight with outer entries $P$ and $-Q$ and a weight with outer entries $c$ and $-c$.
Contracting with the dual of the second weight leads to the flat weight with outer entries $P+2c$ and
$-Q-2c$. On suitable integral elements the composite is multiplication by a fixed polynomial, up to
powers of $p$ of the size of the logarithm of the degrees.

% before: pd-kappa-cloud-position
Moving these elements by a matrix $g$ gives the same identities for the lattice spanned by the
columns of $g$, and for the lattices obtained by scaling its first column. These are the lattices
that occur at the source, at an edge and at the target.

% before: pd-kappa-cloud-maps
Weights with small middle entries are reached from flat weights by multiplying with vectors that
carry the middle part, as many as its dimension. Contracting back lands in the flat weight with
exponents raised by $2c_{cl}$, and the composite is the same map for all of them, up to a power of
$p$ of the size of the logarithm of the degrees.

% before: pd-kappa-cloud-retraction
The map between the two flat weights is controlled by its effect on the lattices $R(e)$. A linear
map that detects membership in a lattice has an integral left inverse, because over a discrete
valuation ring a saturated sublattice is a direct summand.

% before: pd-kappa-cloud-smith
Split families on the two flat weights have to be combined into one family. The combination
reduces a matrix of pairings over a discrete valuation ring, separating the rows that can be
inverted up to $z_1z_2$ from those whose entries are multiples of $z_1z_2$ by elements of the
maximal ideal.

% before: pd-kappa-cloud-upper
The rows of the second kind give a split family whose forms are only known to lie in $R'(e)$. The
bound $n\leq\mu(\lambda_0)$ of Lemma~\ref{pd-kappa-rank-upper} holds for such families too, which
limits the number of these rows.

% before: pd-kappa-cloud-count
With these tools the maps from flat weights carry split families on two flat weights to a split
family on a weight with small middle entries. Its size is the number of maps times a number that
falls short of the size of the first family by at most the defect of the second.

% before: pd-kappa-weight-change
The defects are estimated for flat weights, while the projection property is needed for weights
with small middle entries and small entries outside the top positions. The passage from flat
weights to such weights shifts the exponents by $c_{cl}$ and costs a bounded amount of level,
margin and depth.

% before: pd-kappa-product-space
At a top position the flat weight with entries $\alpha$ and $-\beta$ is the harmonic part of the
polynomials of degree $\alpha$ in one vector and degree $\beta$ in a dual vector. The ranks for flat
weights are obtained from ranks on these product spaces, which contain the flat pieces of all
lower degrees.

% before: pd-kappa-dimensions
The decomposition of a product space into flat pieces can be seen on dimensions. These counts turn a
rank on a product space into a sum of ranks on the flat pieces weighted by their dimensions.

% before: pd-kappa-harmonic-minors
The dual Weyl module of $(P,0,\dots,0)$ only sees the first column, and that of $(0,\dots,0,-Q)$ only
sees the $m\times m$ minors on the first $m$ columns. These minors behave like the coordinates of a
dual vector: they are algebraically independent, they pair to zero with the first column, and a
matrix acts on them through its inverse transpose.

% before: pd-kappa-harmonic-decomposition
Polynomials in a vector and a dual vector decompose into harmonic pieces times powers of the invariant
$h=\sum_ix_iy_i$. The projections onto the pieces are explicit polynomials in multiplication by $h$
and the Laplacian $\Delta$, with binomial coefficients in the denominators.

% before: pd-kappa-harmonic-fock
The denominators of these projections are controlled on a lattice weighted by factorials, which the
divided powers of $h$ and of $\Delta$ preserve. Comparing it with the polynomials with integral
coefficients costs factorials whose $p$-adic valuation is of the size of the logarithm of the degree.

% before: pd-kappa-harmonic-model
The two-letter spaces and the flat dual Weyl modules have explicit models in polynomials in a vector
and a dual vector, in which the flat pieces correspond to the harmonic polynomials. In these models
the lattices become the polynomials with integral coefficients after a change of variables, and the
harmonic decomposition has denominators bounded by a power of $p$.

% before: pd-kappa-harmonic-gauss
The pieces of the decomposition are harmonic polynomials multiplied by powers of $h$. Multiplication
by $h$ neither creates nor removes denominators, which is seen, as in Gauss's lemma, on the monomial
with the largest exponent of $x_1$.

% before: pd-kappa-harmonic-single
At a single top position the product space is a product of two dual Weyl modules, of degree $P$ in
one vector and degree $Q$ in a dual vector. It decomposes into the flat pieces of degrees $P-i$ and
$Q-i$, and the projections onto them have denominators bounded by a power of $p$ of the size of the
logarithm of the degrees.

% before: pd-kappa-harmonic-maps
The decomposition of a product space into flat pieces is realized by projections and inclusions, one
pair for each function $j$. They are tensor products of the maps at the single positions, and they
carry the cores of the product space into the cores of the pieces up to a power of $p$.

% before: pd-kappa-harmonic-block
Projecting a split family on a product space to the pieces gives matrices of pairings whose sum is a
multiple of the identity. Over a discrete valuation ring such a sum splits into pieces whose ranks
add up to at least the size of the family.

% before: pd-kappa-harmonic-split
A split family on a product space projects to split families on its flat pieces. The projections
are integral up to a power of $p$ of the size of the logarithm of the degrees, which is the loss
$4|T|u_0$ in the depth.

% before: pd-kappa-function
The flat defect of a pair of exponents is the least defect for which the edge rank condition
holds, at a level and a depth fixed by the data. The function of the theorem combines the flat
defects at the exponents shifted down and up by $c_{cl}$, which is what the passage to general
weights requires.

% before: pd-kappa-defects-shares
The defects of the segments measure how much of a lattice fails to pair well with the whole lattice.
A saturated sublattice is good when each of its primitive vectors has a partner whose pairing with it
has valuation at most the depth, and the bad share is the proportion of the rank that the largest
good sublattice misses.

% before: pd-kappa-defects-bad-share
The bad share is the defect of a good share. This follows from the definition, because the zero
sublattice is good and the ranks of good sublattices are bounded by the rank of the lattice.

% before: pd-kappa-defects-segment-lattice
Each segment between consecutive even elements of $T$ carries a lattice of polynomials in a chain of
letters, one vector and one covector letter for each odd position and a vector letter for the next even
position. It is a sum over blocks, which record the $p$-adic digits of the exponents at the positions
and the degrees given to the letters, of lattices cut out by conditions on the weights of the
monomials and on the valuations of their coefficients.

% before: pd-kappa-defects-numbers
The margin term and the segment defects are formed from the data of
Theorem~\ref{pd-gain-theorem-one} with a few more constants: one for the error term of the segments, one
for the margin, one for the margin term and one for the depth. The segment datum of a segment records
its gap, the distances of its odd positions, the exponents $b_T$ there and at the next even position, and
how many of the distances are short.

% before: pd-kappa-defects-margin
The margin term is $C_\kappa|T|$ times a sum of four terms. Once the constant $C_H$ is large compared
with the other constants, the hypotheses of Theorem~\ref{pd-gain-theorem-one} bound each of them by a
multiple of $N_1/Y$.

% before: pd-kappa-defects-windows-blocks
The rank bound on the product cores is assembled from the segment lattices block by block. Each block
spans a finite-dimensional space over $\mathcal K$, and the pairs of lattice vectors used are supported
on single blocks, where the segment pairing is a multiple of the Fock pairing.

% before: pd-kappa-defects-windows-ranks
The rank of a segment lattice is the dimension of the space it spans, because every vector of that
space has a nonzero multiple in the lattice. That space is the sum of the spaces of the kept blocks.

% before: pd-kappa-defects-windows-partners
A good sublattice of a segment lattice yields as many pairs of vectors, each supported on one block and
paired diagonally, as its rank. The Fock pairing on each block lattice has bases in elementary divisor
form, and a good sublattice meets the span of the basis vectors whose pairing has large valuation only
in $0$.

% before: pd-kappa-defects-windows-model-segment
The lower bound for the product of the ranks of the segment lattices comes from a lower bound for each
segment separately. For one segment it compares the sum of the dimensions of the kept blocks with a
product of dimensions $P_m$, one for each letter of the segment, in which the exponent of the first
letter is replaced by its coarse value $\hat a_j$.

% before: pd-kappa-defects-windows-model-step
The first letter of a segment carries the exponent $a_T(s_j)$, while the segment bound uses
$\hat a_j$, which is close to $a/p^{g_j}$. Since $a_T(s_j)$ lies in a short window above $a$, the ratio of
$P_m(a_T(s_j))$ to $p^{g_jm}P_m(\hat a_j)$ is close to $1$. The losses in this comparison and in the
factor in front of the segment bound add up to at most three times the bracket of the margin term.

% before: pd-kappa-defects-windows-model
The product of the ranks of the segment lattices accounts for the dimension of the product space up to
the factor $[k]$ and the margin term. This is the counting part of the rank bound.

% before: pd-kappa-defects-windows-lift-pairs-model
The lifts of the segment families are polynomials in one row of variables for each letter $(s,1)$ and $(s,0)$ of the
product space. A flat theta map identifies these polynomials with the product space so that integrality and the
divisibility at each edge become conditions on coefficients. Each segment polynomial is moved into these variables by a
substitution with $p$-power exponents, after multiplication by the monomials that its digits and its top letters fix.

% before: pd-kappa-defects-windows-lift-pairs-flat-dual
The dual of an edge is again an edge, and a basis matrix for it is the inverse transpose of one for the edge, with its
first column multiplied by $p$. So the flat theta maps of an edge and of its dual edge are attached to two matrices that are
inverse transposes of each other up to that correction.

% before: pd-kappa-defects-windows-lift-pairs-flat-theta
Flat theta maps exist for every invertible matrix. On a vector letter the map takes the first column of $h_s^{-1}x_{s,1}$,
and on a covector letter the last row of the adjugate of $x_{s,0}$ times $h_s$, as for the flat models of
Lemma~\ref{pd-symbols-flat-exists}.

% before: pd-kappa-defects-windows-lift-pairs-flat-forms
Read through the flat theta map at the identity matrix, the Fock pairing scaled by $z_T$ becomes a map from the product
space to its dual. The scaled Fock pairing of two integral polynomials of the given row degrees is integral, and a change
of variables on one side is undone by the inverse transpose on the other. So this map carries the lattice of a dual lattice
into the dual of the lattice, and the module of a dual edge into the dual module of the edge.

% before: pd-kappa-defects-windows-lift-pairs-flat
Two flat theta maps are needed, one attached to the edge and one to its dual edge, and a map from the space to its dual
carries the cores of the dual edge to the dual cores of the edge. Through these maps the edge pairing of the product space
becomes the Fock pairing of the substituted polynomials, scaled by $z_T$.

% before: pd-kappa-defects-windows-lift-pairs-source-symbols
The symbol families are integral, and their factors for different segments involve different letters. A kept block can be
read off from any monomial of its image, because the digits of the block are the residues of the exponents modulo the
$p$-powers $P_i$. So images of distinct kept blocks are orthogonal for the Fock pairing, and the images of kept blocks have
the row degrees of the product space.

% before: pd-kappa-defects-windows-lift-pairs-source-order-segment
The exponent of a kept block is bounded one segment at a time, from a few facts about the segment: the windows for the
degrees $b_i$ around $b'$, the gap between $a/b'$ and the odd powers of $p$, and the spacing of the odd distances $d_i$. For
each letter the gap decides which term of the minimum is the smaller one, and the conditions on a kept block bound the
letter's contribution by that term, up to errors of size $b'/p$, $\Lambda a/(pk)$ and $\Lambda b'/k$. The same estimates
bound the cap of the block, and with it the entry $o_+$.

% before: pd-kappa-defects-windows-lift-pairs-source-order
The lift of a symbol family is integral only after multiplication by a power of $\varpi_0$ that grows with the blocks. The
exponent is bounded segment by segment, and the sum of these bounds is at most $t_0$, which is the depth allowed in the
cores.

% before: pd-kappa-defects-windows-lift-pairs-source-lift-relaxed
The lift is built by successive approximation, and each step produces symbol families of blocks that are no longer kept.
The blocks that occur satisfy weaker bounds with a slack, and the families are collected in modules indexed by the margin,
the slack and the order.

% before: pd-kappa-defects-windows-lift-pairs-source-lift-kept
The approximation starts from the kept blocks. The bounds of a kept block imply those of a relaxed block of slack $0$,
by the windows for $b_T$ and the band condition on $a/b'$.

% before: pd-kappa-defects-windows-lift-pairs-source-lift-tail-systems
One step of the approximation can be taken on each symbol family separately, as the families it applies to form a module.
A family is lifted by writing it as a product of leading terms of polynomials that are already lifted. The products of
these leading terms with monomials of nearby degrees must then fall into the symbol modules of higher order.

% before: pd-kappa-defects-windows-lift-pairs-source-lift-tail-divisible
The substitution at the identity multiplies each variable by $1$ or by $\varpi_0$. The substitution at any matrix of
$GL_k(W)$ keeps the same divisibility on the coefficients of an integral polynomial, which bounds the error terms of a
factorization.

% before: pd-kappa-defects-windows-lift-pairs-source-lift-tail-factorization
A factorization system determines its lift, the product of its factors divided by the power of $\varpi_0$ given by the
order. Expanding the substitution of this product around the leading terms gives the family itself, terms divisible by
$\varpi_0^n$, and terms in the symbol modules of higher order.

% before: pd-kappa-defects-windows-lift-pairs-source-lift-tail-columns
The polynomials $f$ of the generators are combinations of path generators, by the path basis of the letter spaces over
$\mathcal R$, and the polynomials $g$ are combinations of monomials. So only symbol families whose polynomials are single
path generators and single monomials need to be lifted.

% before: pd-kappa-defects-windows-lift-pairs-source-lift-tail-existence-situation
A factorization system for a symbol family of column generators is assembled from three constructions. They share the data
of the lemma that assembles them: the relaxed blocks, the margin $\mu+n\kappa_1$ of the generators and the order $o$.

% before: pd-kappa-defects-windows-lift-pairs-source-lift-tail-existence-levels
A column generator is determined by a level of its block, a path generator and a monomial in the $\xi_c$. Its image under
$\iota_j(\beta_j,\cdot)$ is the image of the factors $\hat E_l$, $\hat E_+$ and the path generator, times a monomial that does not
depend on the path generator. The factors of a factorization system are lifted factors: their substitution at the identity is a
power of $\varpi_0$ times a given polynomial, up to a higher power of $\varpi_0$.

% before: pd-kappa-defects-windows-lift-pairs-source-lift-tail-existence-column-systems
A path generator is a product of polynomials of degree at most one in each letter. A column system records such a
factorization together with a lifted factor of the image of each factor times a monomial. The monomials add up to the part of
the base monomial at the first variables of the letters $(s_j,1)$ and $(y_{j,i},1)$.

% before: pd-kappa-defects-windows-lift-pairs-source-lift-tail-existence-columns-data
The column system of a path generator comes from a sequence of rank $k$ that extends the path by a first entry. That entry
counts, letter by letter, the exponents at the first variables that the base monomial provides, and an extra first letter at
$s_j$ carries the remaining power of $z_{(s_j,1),1}$. Its lifted column functions are the lifted factors.

% before: pd-kappa-defects-windows-lift-pairs-source-lift-tail-existence-columns-numerics
The first entry of the augmented sequence starts at $a_0$, which must be nonnegative, and it must bound the first entries of
the path at the vector letters $x_i$ with $i\leq\eta$. Both are numerical bounds on the orders $\varsigma_i$ and $w_i$ of a
relaxed block.

% before: pd-kappa-defects-windows-lift-pairs-source-lift-tail-existence-columns-shape
The lifted column functions of the augmented sequence behave like those of a path once its first entry is nondecreasing,
constant across the covector letters, and above the other entries at the vector letters. At the letters $x_{u_l}$ and $x_+$
this is the bound on the path that defines the column generators; at the $x_i$ with $i\leq\eta$ it is the numerical bound
above.

% before: pd-kappa-defects-windows-lift-pairs-source-lift-tail-existence-columns-lifted
The augmented sequence is not a path, and its plane part has rank $1$ only. The bounds on the substitutions of a lifted column
function need of the sequence only how the numbers of entries at least $c$ change from letter to letter, and the congruence needs
in addition that its vectors be dominant; a path with $N+r\leq K$ has these properties.

% before: pd-kappa-defects-windows-lift-pairs-source-lift-tail-existence-columns-augmented
The shape of the augmented sequence gives the properties needed for its lifted column functions. A rise in the number of entries
at least $c$ of the path at a vector letter forces $c$ below the first entry there, so the first entry rises with it; and only
vector letters can be costly, at positions beyond those of the earlier letters.

% before: pd-kappa-defects-windows-lift-pairs-source-lift-tail-existence-columns-symbol
The lowest part of a lifted column function of the augmented sequence splits into a plane factor and a tangent factor. After the
renaming $\phi$, the first is a monomial in the first variables of the letters, and the second is the image of a column function
of the path.

% before: pd-kappa-defects-windows-lift-pairs-source-lift-tail-existence-columns-counts
The column data, summed or multiplied over the columns, give back the data of the path generator. The degrees and the column
functions recover those of the path, and the monomials and the orders add to the exponents and the order that a column system
prescribes.

% before: pd-kappa-defects-windows-lift-pairs-source-lift-tail-existence-columns-degrees
The remaining conditions of a column system concern the degrees of the column data. They come from the degrees of the column
functions and of the substitutions of the letters.

% before: pd-kappa-defects-windows-lift-pairs-source-lift-tail-existence-columns-clauses
With the augmented sequence in this shape, its column functions give the column system. The path generator is the product of
the column functions of the path, and the substitution at the identity of each lifted column function is, up to a higher power
of $\varpi_0$, a power of $\varpi_0$ times the image of a column function times a monomial in the first variables.

% before: pd-kappa-defects-windows-lift-pairs-source-lift-tail-existence-columns
Every path generator that occurs in a column generator has a column system. Its factors are column functions of the path, whose
product is the path generator.

% before: pd-kappa-defects-windows-lift-pairs-source-lift-tail-existence-prefactorizations
Besides the column factors, a factorization system for a symbol family needs lifted factors of the images of the $\hat E_l$ and
of monomials, and a factor whose substitution at the identity is exact. A prefactorization collects these with the identities of
row degrees, orders and exponents that make their product a lift of the family. The tail condition is the last condition of a
factorization system, written for the factors of a prefactorization.

% before: pd-kappa-defects-windows-lift-pairs-source-lift-tail-existence-prefactorization-keys
The factors of a prefactorization beyond the column factors are built from a few polynomials. The images of the $\hat E_l$
and of $\hat E_+$ are sums of products of the variables of the two letters at one position. The other polynomials lift the
monomials of the digits and of the orders $\varsigma_i$, and the polynomial $L_{j_0}$ carries the bound at the planes other
than that of $1$.

% before: pd-kappa-defects-windows-lift-pairs-source-lift-tail-existence-prefactorization-products
The factors of a prefactorization are built as products and powers of a few lifted factors. Products of lifted factors are
lifted factors of the products of the polynomials they lift, with the orders added. A lifted factor can be read off from an
exact formula for its substitution at the identity.

% before: pd-kappa-defects-windows-lift-pairs-source-lift-tail-existence-prefactorization-hooks
The images of the $\hat E_l$ and of $\hat E_+$ are lifted by powers of the polynomials $h_x$. Every substitution
$\Sigma_g$ multiplies $h_x$ by $\varpi_0$, as the vector and the covector letter at $x$ are substituted by inverse
matrices. A $p$-power of $h_x$ agrees with the corresponding $\hat E_{x,P}$ modulo $p$.

% before: pd-kappa-defects-windows-lift-pairs-source-lift-tail-existence-prefactorization-key-factors
The monomials of the digits and of the orders $\varsigma_i$ are lifted by the polynomials $\kappa^S_{j,i}$ and
$\kappa^W_{j,i,t}$, and a power of the variable $z_{(s_j,1),1}$ by the $D_t$. Modulo $\varpi_0$ their substitutions involve
only the first row of the substituted matrix at $s_j$, and its $p^d$-th power is the first row at $y_{j,i}$, by the
Frobenius congruence.

% before: pd-kappa-defects-windows-lift-pairs-source-lift-tail-existence-prefactorization-localizer
The product of a prefactorization must have substitutions divisible by $\varpi_0^{o+c}$ at the planes other than that of
$1$, while the orders of its factors add up to $o$. The remaining power $\varpi_0^c$ comes from $L_{j_0}^c$: modulo
$\varpi_0$ its substitution is a multiple of a difference of two products that the Frobenius congruence makes equal.

% before: pd-kappa-defects-windows-lift-pairs-source-lift-tail-existence-prefactorization-room
The monomial factor of a prefactorization is what remains of the base monomial after the other factors are removed. It
has a nonnegative exponent at the first variable of each letter $(s_{j+1},1)$ only if the exponents used there by the other
factors are at most $a_T(s_{j+1})-o_+$. This is a numerical bound on the digits, on the orders $\varsigma_i$ and on the error
exponent $E_0$.

% before: pd-kappa-defects-windows-lift-pairs-source-lift-tail-existence-prefactorization-family
The identities of row degrees and orders in a prefactorization are read off from the symbol family. Its factor for the
$j$-th segment is homogeneous for every balanced weight, and two of these weights give the row degrees and the
$\varpi_0$-degree.

% before: pd-kappa-defects-windows-lift-pairs-source-lift-tail-existence-prefactorization
The column systems of the path generators of a symbol family extend to a prefactorization. The bound on the substitutions of
the product at the planes other than that of $1$ comes from a further factor with substitutions divisible by $\varpi_0^c$ there,
and from the bounds on the substitutions of the single factors.

% before: pd-kappa-defects-windows-lift-pairs-source-lift-tail-existence-segment-numbers
The tail condition replaces a block by a nearby block of a larger slack, for the smaller margin $\mu+(n-e)\kappa_1$. The slack
grows by $e\kappa_1$, which exceeds by a multiple of $N_1$ the changes in the orders and in the numbers $\hat b_i$; so the nearby
block is again relaxed, and its capacities exceed those of the first block by more than the loss in the conditions $Y$.

% before: pd-kappa-defects-windows-lift-pairs-source-lift-tail-existence-segment-block
A monomial in the variables of the letters of one segment determines a block. The residues of its exponents modulo $P_i$
are the digits, and the quotients, with the degrees of a polynomial in the letters, give the orders.

% before: pd-kappa-defects-windows-lift-pairs-source-lift-tail-existence-segment-block-numbers
The numbers of the block of a monomial follow from division with remainder. When the row degrees are the right ones, its
digits are those of a relaxed block, and bounds on the exponents bound the quotients.

% before: pd-kappa-defects-windows-lift-pairs-source-lift-tail-existence-segment-weights
The conditions $Y(\mathcal L,M)$ behave well under products, as the operators $E^{\mathcal L}_{cd}$ are derivations and the
weight is additive. A single variable and a sum $\sum_cz_{a,c}z_{a',c}$ over a vector and a covector letter raise the bound by
at most one.

% before: pd-kappa-defects-windows-lift-pairs-source-lift-tail-existence-segment-quotient
A product of powers of $z_l$ and $z'_l$ is a combination of products of powers of $\hat E_l$ and of the sums
$\sum_cx_{u_l,c}y_{u_l,c}$, by the binomial formula. So a product of the shape of a generator, with $z'_l$ in place of some of
the factors $\hat E_l$, lies in the span of the generators of the block when the remaining factors satisfy the conditions of
the higher levels.

% before: pd-kappa-defects-windows-lift-pairs-source-lift-tail-existence-tail-condition-segment
The tail condition is checked one segment at a time. A product of factors $\hat E_l$, a polynomial in the letters and a
monomial whose exponents are near those of the base monomial is the image of a single block: the residues of its exponents
modulo $P_i$ are the digits of that block, and the quotients add to its orders.

% before: pd-kappa-defects-windows-lift-pairs-source-lift-tail-existence-tail-condition
Removing some factors from a prefactorization leaves the image of a product of factors $\hat E_l$ and of a divisor of the path
generators, times a monomial near the base monomial. The tail condition asks that such products of the right row degrees and
order be combinations of symbol families with a larger slack and a higher order.

% before: pd-kappa-defects-windows-lift-pairs-source-lift-tail-existence-quotient
Removing column factors from a path generator weakens the conditions $Y(\mathcal L,M)$ that it satisfies by a bounded amount.
The weight of a monomial for $\mathcal L$ is at least minus its degree in the covector letters, and the largest weight $\lambda$
is additive on products.

% before: pd-kappa-defects-windows-lift-pairs-source-lift-tail-existence-assembly
A prefactorization becomes a factorization system once its factors are listed with their multiplicities. The last condition of
a factorization system reduces to the tail condition: the bounds on the removed factors come from the conditions of the
prefactorization and from Lemma~\ref{pd-kappa-defects-windows-lift-pairs-source-lift-tail-existence-quotient}.

% before: pd-kappa-defects-windows-lift-pairs-source-lift-tail-existence
For a symbol family of column generators the factors come from the column functions of its path generators, from the
factors $\hat E_l$ of the levels and from monomials. The last condition of a factorization system asks that, when some
factors are removed, the product of the remaining leading terms with a monomial be a combination of symbol families of
relaxed blocks with a larger slack and a higher order.

% before: pd-kappa-defects-windows-lift-pairs-source-lift-tail
Each step of the approximation lifts one symbol family modulo a higher power of $\varpi_0$. The error is a combination of
symbol families of higher order, with a smaller margin and a larger slack, multiplied by powers of $\varpi_0$.

% before: pd-kappa-defects-windows-lift-pairs-source-lift
The lift itself is a polynomial of the row degrees of the product space. Its substitution at the identity matches the symbol
family modulo $\varpi_0^c$, and its substitutions at the planes other than that of the identity are divisible by
$\varpi_0^c$, which is what the cores at the other edges require.

% before: pd-kappa-defects-windows-lift-pairs-source
A symbol family built from kept blocks is matched, modulo $\varpi_0^c$, by the substitution of a polynomial whose flat theta
image lies in the core at the edge. This is how the families of the segments reach the cores of the product space.

% before: pd-kappa-defects-windows-lift-pairs-form-numbers
The scale $z_T$ splits into one factor for each segment, as the letters do. The scaled Fock pairing of integral polynomials
of the given row degrees is integral, and the valuations of the segment factors add up to at most $4\kappa_0$.

% before: pd-kappa-defects-windows-lift-pairs-form-block-data
The block form is built from a differential operator on the polynomials of a segment that acts on each monomial by a
scalar. The scalars are compared with ratios of factorials of the exponents that the map $\iota_j(b,\cdot)$ produces, and
only exponents of fixed weighted degrees occur; all of these depend on the block.

% before: pd-kappa-defects-windows-lift-pairs-form-block-monomials
The map $\iota_j(b,\cdot)$ sends distinct monomials to distinct monomials, so the Fock pairing of two images is a sum over
the monomials of $\mathcal P_j$. The operators act on monomials by scalars, and the factorials of the image exponents split
variable by variable.

% before: pd-kappa-defects-windows-lift-pairs-form-block-homogeneous
The generators of a block have one weighted degree for every balanced weight. Their monomials therefore satisfy a fixed
set of linear conditions, which bound their factorials by the number $N_b$ and fix the factorial powers of the levels.

% before: pd-kappa-defects-windows-lift-pairs-form-block-saturated
The maps of the block form are scalings of the block lattice, which lies in the $\mathcal K$-span of the symbol module.
An element of that span with integral coefficients already lies in the symbol module, by the basis of the block rank.

% before: pd-kappa-defects-windows-lift-pairs-form-block-shift-operators
The operator $X^\tau D_\tau$ of the block form is a composite of operators in one variable each, and each of these
multiplies a monomial by a binomial coefficient of its exponent in that variable. A Leibniz rule computes them on products,
and the generators of a symbol module are products of powers.

% before: pd-kappa-defects-windows-lift-pairs-form-block-shift-spaces
The polynomial in the letters that enters a generator lies in spaces $Y(\mathcal L,M)$ whose bounds carry the margin. A
one-variable operator of degree $n$ keeps it in these spaces once the bounds are raised by $n$, and multiplying it by a
variable of a letter and a variable of a covector letter costs one unit.

% before: pd-kappa-defects-windows-lift-pairs-form-block-shift-generators
An odd variable occurs in at most one of the $\hat E_l$ and otherwise only in the polynomial in the letters or in a power
of $z_l$. Where the operator falls on a power of $\hat E_l$, it trades factors of $\hat E_l$ for products of the two letters
that $\hat E_l$ pairs, which lowers the level and costs one unit of margin for each factor.

% before: pd-kappa-defects-windows-lift-pairs-form-block-shift
The differential operator of the block form must keep the symbol module, at the cost of a smaller margin. Each factor
$X^\tau D_\tau$ at the odd variables lowers the margin by the degree of $\tau$.

% before: pd-kappa-defects-windows-lift-pairs-form-block-integral
The coefficients of an element of the block lattice are integral only after multiplication by a power of $\varpi_0$ that
the factorials of the exponents determine. A pair of such elements has integral products of coefficients scaled by these
factorials.

% before: pd-kappa-defects-windows-lift-pairs-form-block-valuations
The scalar of the block form compares the factorials $\varphi_j$ of the letter degrees of the segment with the factorials
that the monomials of the block produce. Its valuation is bounded by comparing the two at the leading exponent of the
block.

% before: pd-kappa-defects-windows-lift-pairs-form-block-factorial
The factorial ratios of the image exponents are $p$-adic units times a fixed scalar, and their inverses agree modulo
$p^c$ with polynomials in binomial coefficients. These polynomials are the eigenvalues of a differential operator, which
corrects the ratio up to $1+\varpi_0^ch$.

% before: pd-kappa-defects-windows-lift-pairs-form-block
The pairing is computed one segment and one kept block at a time. For a kept block, two maps send the block lattice into
the symbol module so that the scaled Fock pairing of the images is a fixed scalar times the block pairing, modulo
$\varpi_0^c$.

% before: pd-kappa-defects-windows-lift-pairs-form
The symbols attached to families carried by tuples of kept blocks have to pair correctly. Any two polynomials close to
their symbol families have Fock pairing, scaled by $z_T$, equal modulo $\varpi_0^c$ to a scalar of bounded valuation times
the product of the segment pairings, and to $0$ when the tuples differ.

% before: pd-kappa-defects-windows-lift-pairs
The lifts are built for one tuple of blocks at a time, one block in each segment. For families carried
by such tuples, the pairing of the lifts is, modulo $\varpi_0^c$, the product of the segment pairings
times a scalar $\kappa_b$ of valuation at most $4\kappa_0$, and it vanishes when the tuples differ.

% before: pd-kappa-defects-windows-lift
The diagonal pairs of the segments are lifted to vectors in the product cores. Their pairings there
are diagonal modulo $\varpi_0^c$, with diagonal entries of valuation at most the sum of the depths plus
$4\kappa_0$.

% before: pd-kappa-defects-windows-perturb
A pairing matrix that is diagonal modulo $\varpi_0^c$, with diagonal entries of valuation less than
$c/4$, can be corrected by a matrix over $\mathcal R$ into a scalar multiple of the identity. The scalar
is the diagonal entry of largest valuation, so the depth does not grow.

% before: pd-kappa-defects-windows
The rank bound on the product cores combines the good shares of the segment lattices. The depth adds
up over the segments, and the defect is the margin term plus the sum of the segment defects.

% before: pd-kappa-defects-segment-mean-operators
The bound for the mean segment defect rests on conditions about operators of $\mathfrak{gl}_n$ type,
their weights and isotypic projectors, and about functions on $GL_{r+1}$ that transform by a character
of the lower triangular matrices. The lattice $\Delta(v)$ is generated by $v$ under the divided powers,
and $\nabla(\lambda,v)$ consists of the vectors whose images under the divided powers have an integral
multiple of $v$ as their component of weight $\lambda$.

% before: pd-kappa-defects-segment-mean-conditions
The mean bound holds for an instance once eleven conditions hold for it. Conditions (1) to (3) concern
the operators on the letters of the kept blocks of the segment, (4) to (7) the functions on $GL_{r+1}$
that transform by a character, (8) the integral polynomials of bounded weight, and (9) to (11) the
ranks of the segment lattices and of their parts with given digits.

% before: pd-kappa-defects-segment-mean-theorem-a-classes
The mean over the window is computed by grouping the kept blocks according to the residues of their
vectors modulo the powers $P_i$. A digit class fixes these residues, and condition (11) compares the
numbers $d(\delta;\alpha)$ with the ranks of the blocks of the class. Their sum $M(\alpha)$ over the classes
is a product of binomial coefficients, one for each entry of $\alpha$.

% before: pd-kappa-defects-segment-mean-theorem-a-average
A mean of quotients of sums is bounded by exchanging the summation over the window with the summation
over the digit classes. The exchange needs two inputs: the rank of each class is at least half of its
number $d$, and the totals of the numbers $d$ change by a bounded factor across the window.

% before: pd-kappa-defects-segment-mean-theorem-a-flat
The averaging needs the total $M(\alpha)$ to vary little with $\alpha$. It is a constant times a product of
binomial coefficients in the entries of $\alpha$, and the window has length about $a/(kY)$, so the ratio
of two totals is at most $e^3$.

% before: pd-kappa-defects-segment-mean-theorem-a-blocks
The segment lattice is the direct sum of the lattices of its kept blocks, and the segment pairing is the
sum of the block pairings. So good sublattices of the blocks combine into a good sublattice of the
segment, and ranks add. A good share bounds the segment defect, which is the smallest defect of a good
share.

% before: pd-kappa-defects-segment-mean-theorem-a-levels
The bound for the defect is assembled level by level. A level fixes the powers of $\hat E_l$ and $\hat E_+$
in a generator of a block lattice; the remaining factors are a polynomial in the letters and a polynomial
in the $\xi_c$, and their products span the level lattice. The badness of a level is the part of its rank
that dual families with pairings $p^D$ do not reach, for a pairing built from the Fock pairing and the
operators $\phi$ attached to the ties.

% before: pd-kappa-defects-segment-mean-theorem-a-free
Moving $\alpha$ within the tuples compatible with a digit class does not change which blocks of that
class are kept, nor their levels. The entries of $\alpha$ enter these conditions only through bounds by
the numbers $\hat\alpha_i$, which are close to $a/P_i$ and larger than the quantities they bound.

% before: pd-kappa-defects-segment-mean-theorem-a-ratio
The mean value estimate for the polynomial parts compares a tuple with the tuples obtained by raising one
entry to the top of the window. Raising an entry changes one binomial factor of $d(\delta;\alpha)$, and as
the window is short, by a factor at most $e^3$.

% before: pd-kappa-defects-segment-mean-theorem-a-depth
The levels are compared at depth $YD_1(C_M)$, and passing from the levels to the block adds
$D_{\mathrm{lv}}$. Both fit within the depth $D_A$ of the good shares once $C_A$ is large, by a computation
with $s\leq N_1$ and $Y\geq1$.

% before: pd-kappa-defects-segment-mean-theorem-a-block-bound-base
The levels are studied over $\mathbf Q_p$, while the block lattices live over $\mathcal R$. The ring $\mathcal R$ is
complete with $p$ in its maximal ideal, so $\mathbf Q_p$ embeds in $\mathcal K$, and applying the embedding to
coefficients carries generators over $\mathbf Q_p$ to generators over $\mathcal R$, keeping linear independence
and the Fock pairing.

% before: pd-kappa-defects-segment-mean-theorem-a-block-bound-rank
The rank of a block lattice is the dimension of the span of its generators, which is the sum of the spans
of the levels. Each of these is spanned by images of the level lattice over $\mathbf Z_p$, since every
polynomial part over $\mathcal K$ is a combination of images of polynomial parts over $\mathbf Z_p$.

% before: pd-kappa-defects-segment-mean-theorem-a-block-bound-dual
A dual family over $\mathbf Z_p$ whose pairings are one power $p^e$ gives a good sublattice of the block
lattice. Scaled by a power of $\varpi_0$ fitted to $e$, the family lies in the block lattice, and the
sublattice it spans in the block lattice is good at the depth of $p^e$.

% before: pd-kappa-defects-segment-mean-theorem-a-block-bound-families-tie
A level lattice is carried into the span of the generators by the operators of its level, after a projection
at each tie. The projection corrects the product of a polynomial part with the powers of the $z_l$ so that the
Fock pairing separates the levels and, on one level, becomes the level pairing up to the integer $c_\omega$.

% before: pd-kappa-defects-segment-mean-theorem-a-block-bound-families-fock-raising
The pairings of tie projections are computed one tie at a time. Multiplication by $\sum_cw_ce_c$ raises both
degrees by one and is adjoint, for the Fock pairing, to the operator $F$; on polynomials that $F$ kills, $F$ undoes
a power of the multiplication up to an explicit product of integers, so different powers are orthogonal.

% before: pd-kappa-defects-segment-mean-theorem-a-block-bound-families-fock-tie
The tie projection $\pi$ is built so that its image is killed by the operator $\hat F$ that includes the pair
$(z,z')$. The previous lemma then applies with $\hat E$ and $\hat F$, and the pairing of two projections reduces to
the pairing of the original polynomials against $\phi$, times $\zeta!$ from the power of $z$.

% before: pd-kappa-defects-segment-mean-theorem-a-block-bound-families-fock
On the images of the level lattices, the Fock pairing is block diagonal. On one level it is $c_\omega$ times the
pairing of the level lattice with the operator $\Phi_\omega$, which is $c_\omega d_\omega$ times the level pairing.

% before: pd-kappa-defects-segment-mean-theorem-a-block-bound-families-lattice-lowering-weights
The level spaces are cut out by the spaces $Y(\mathcal L,M)$ of Definition~\ref{pd-kappa-defects-segment-lattice}(4),
the largest subspaces of a span of monomials of bounded weight that are stable under the operators
$E^{\mathcal L}_{cd}$. An operator that commutes with these operators and does not raise weights preserves such a
space. Derivatives in a vector letter, and the operator of a tie whose two letters lie in the segment, are of this
kind.

% before: pd-kappa-defects-segment-mean-theorem-a-block-bound-families-lattice-lowering-pairs
Raising one entry of a level need not give a level. With the degrees lowered with truncation, the spaces and
lattices of a level are defined for every pair of tuples, and the lowering is stated for all pairs.

% before: pd-kappa-defects-segment-mean-theorem-a-block-bound-families-lattice-lowering-spaces
Raising $j_l$ by one lowers the degrees of $x_{u_l}$ and $y_{u_l}$ by one, as the operator $F_l$ of the tie does, and
raising $j_+$ lowers the degree of $x_+$, as a derivative in a variable of $x_+$ does. By the previous lemma both
operators also preserve the conditions of the initial segments.

% before: pd-kappa-defects-segment-mean-theorem-a-block-bound-families-lattice-lowering-sum
At the tie of $x_+$ and $\xi$ the operator $F_+$ acts on both factors of the products that span a lattice. It is a
sum of $m$ commuting operators, one for each coordinate, and a power of the sum is expanded by the binomial theorem,
at the cost of the inverses of binomial coefficients, whose valuation is at most $L$.

% before: pd-kappa-defects-segment-mean-theorem-a-block-bound-families-lattice-lowering
Undoing a tie projection produces powers of the operator $F_a$ of the tie applied to a level lattice. These
powers, divided by the squares of factorials and multiplied by a fixed power of $p$, carry the lattice of a pair
$(j,j_+)$ into the lattice of the pair with $j_a$ raised.

% before: pd-kappa-defects-segment-mean-theorem-a-block-bound-families-lattice-step
The tie projections are removed one tie at a time, starting with the tie at $x_+$. Expanding one projection writes
it as a combination of the unprojected products for higher levels, with coefficients whose denominators are
binomial coefficients of bounded valuation, so each removal costs a factor $p^u$.

% before: pd-kappa-defects-segment-mean-theorem-a-block-bound-families-lattice
The images of the level lattices have coefficients in $\mathbf Q_p$ with bounded denominators. A power $p^U$, with
$U$ of the size of the number of ties times a logarithm of the degrees, carries them into the span of the
generators over $\mathbf Z_p$.

% before: pd-kappa-defects-segment-mean-theorem-a-block-bound-families-valuation-factorials
The valuation of the scalars of a level is compared with the factorials of a monomial through binomial
coefficients. A binomial coefficient with top entry below $p^{L+1}$ has valuation at most $L$, so splitting a
factorial into the factorials of the parts of a sum loses at most $L$ for each part.

% before: pd-kappa-defects-segment-mean-theorem-a-block-bound-families-valuation-degrees
The factorials of a monomial are compared with those of the degrees of the generators. Grouping each
variable $z_l$, $z'_l$ with a letter of its tie makes every generator homogeneous in each group, and the
powers of $\hat E_l$ and $\hat E_+$ then drop out of the degrees, which are the same at every level.

% before: pd-kappa-defects-segment-mean-theorem-a-block-bound-families-valuation
The scalars $c_\omega d_\omega$ and the power $p^{2U}$ enter the pairings of the assembled family. Their valuation
is bounded by the factorials of the monomials of the generators, up to the loss $D_{\mathrm{lv}}$.

% before: pd-kappa-defects-segment-mean-theorem-a-block-bound-families
The dual family is assembled from the dual families of the level pairings. Its size falls short of the
sum of the ranks of the level lattices by at most the badness of the levels, and the valuation of its
pairing is bounded by the factorials of the monomials up to the loss $D_{\mathrm{lv}}$.

% before: pd-kappa-defects-segment-mean-theorem-a-block-bound
A block lattice has a good sublattice whose rank falls short of the rank of the block by at most the
badness of its levels. The levels contribute dual families, and the depth grows by $D_{\mathrm{lv}}$ in
passing from the levels to the block.

% before: pd-kappa-defects-segment-mean-theorem-a-plus-pairing
Only the factor of $\Phi_\omega$ at $x_+$ depends on $\alpha_+$. Removing it gives a pairing on the level lattice that
is the same for the whole window of $\alpha_+$, and its dual families can be chosen once for all $\alpha_+$.

% before: pd-kappa-defects-segment-mean-theorem-a-plus-rank
The level lattice is spanned by products of polynomial parts with monomials in the $\xi_c$. Its rank is
therefore at most the dimension of the polynomial part times the number of these monomials.

% before: pd-kappa-defects-segment-mean-theorem-a-plus-phi
The pairing without the factor at $x_+$ does not involve the $\xi_c$, so it splits along the monomials in the
$\xi_c$. A dual family of the polynomial part, multiplied by each monomial, gives a dual family of the level
lattice, and the badness grows at most by the number of monomials.

% before: pd-kappa-defects-segment-mean-theorem-a-plus-lattice
A symmetric pairing on a lattice over $\mathbf Z_p$, with $p$ odd, has an orthogonal basis. The basis vectors whose
pairing with themselves has small valuation form a family that is as large as any dual family of that depth.

% before: pd-kappa-defects-segment-mean-theorem-a-plus-orth
The pairing without the factor at $x_+$ is symmetric, as the operators of the ties are self-adjoint for the
Fock pairing and commute. So the level lattice has an orthogonal family whose size misses its rank by at
most the badness of this pairing.

% before: pd-kappa-defects-segment-mean-theorem-a-plus-orth-general
The lemma is stated for the levels and the depth at which it is used. Its proof uses neither, since the
symmetry of the pairing and the finiteness of the level lattice hold for every level of every block.

% before: pd-kappa-defects-segment-mean-theorem-a-plus-window-level-bounds
The window bound needs the pairings of a level to be integral after a fixed power of $p$. The operators $\phi$
divide by binomial coefficients, and a binomial coefficient below $p^{L+1}$ has valuation at most $L$, so each
of these operators enlarges the coefficients of a polynomial by at most a fixed power of $p$.

% before: pd-kappa-defects-segment-mean-theorem-a-plus-window-level-bounds-sharp
Part (2) of the lemma passes through the integrality of a multiple of $\phi(w,e;o,n,\zeta)u$, which costs a
further factor $p^{r+1}$. Its proof gives the bound before that step, for every bound on $u$.

% before: pd-kappa-defects-segment-mean-theorem-a-plus-window-level-pair
The level pairing depends on the exponent $\zeta$ of $z_+$ only through the factor $\phi$ at $x_+$. After
multiplication by one binomial coefficient in $\zeta$, the pairing becomes a combination of the binomial
coefficients $\binom\zeta i$, $i\leq J$, with coefficients that do not depend on $\zeta$.

% before: pd-kappa-defects-segment-mean-theorem-a-plus-window-level-lead
The Gram matrices of the level pairings are therefore the values of one matrix of polynomials in $\zeta$. When its
top coefficients form a diagonal matrix, the leading term of its determinant is their product, divided by a power
of $J!$.

% before: pd-kappa-defects-segment-mean-theorem-a-plus-window-level-smith-gauss
The count of small elementary divisors goes through Gauss norms of polynomials over a field with an ultrametric
absolute value. These norms are multiplicative, and the Gauss norm of the determinant of a matrix of polynomials is
at most the product of bounds on its rows.

% before: pd-kappa-defects-segment-mean-theorem-a-plus-window-level-smith-binomial
The integers congruent to $t$ modulo $p^h$ are the values of $p^hX+t$ at the $p$-adic integers. Composed with this
map, a binomial polynomial $\binom Xk$ with $k\leq p^h$ still has integral coefficients, so the matrix $M$ composed
with it has entries of Gauss norm at most $1$.

% before: pd-kappa-defects-segment-mean-theorem-a-plus-window-level-smith-point
At one integer $t$, each small elementary divisor of $M(t)$ makes the determinant of $M$ small on the disc of
radius $p^{-h}$ around $t$. Comparing with the factorization of the determinant bounds their number by the
distances from $t$ to its roots, measured down to $p^{-h}$, and by the number of roots in that disc.

% before: pd-kappa-defects-segment-mean-theorem-a-plus-window-level-smith-count
The distances from the integers of a progression to one point are controlled by divisibility, since two
integers close to the same point are congruent modulo a power of $p$. When the progression has at most $p^h$
terms, a root lies within $p^{-h}$ of at most one of them.

% before: pd-kappa-defects-segment-mean-theorem-a-plus-window-level-smith-mean
The bound at one integer is summed over a window. With the two counts of the previous lemma, this bounds the mean
number of small elementary divisors by the degree and the leading coefficient of the determinant.

% before: pd-kappa-defects-segment-mean-theorem-a-plus-window-level-smith
At each integer $\zeta$ the badness of the pairing is bounded by the number of small elementary divisors of the
value of the matrix. Their count, summed over a window of consecutive integers, is controlled by the roots of
the determinant and the size of its leading coefficient.

% before: pd-kappa-defects-segment-mean-theorem-a-plus-window-level-gram
A dual family for a bilinear form can be read off an elementary divisor decomposition of its Gram matrix. The
elementary divisors of valuation at most $D'$ give pairings $p^{D'}$, so the badness at $D'$ is at most the number
of the others, and the count of the previous lemma bounds its sum over the window.

% before: pd-kappa-defects-segment-mean-theorem-a-plus-window-level
For one level, the factor at $x_+$ depends on the exponent of $z_+$, and consecutive exponents can be compared.
Summed over $W'$ of them, the badness on the span of an orthogonal family is bounded by its size times a
share of order $J/W'$, plus a term that decreases with the depth $D'$.

% before: pd-kappa-defects-segment-mean-theorem-a-plus-window-numbers
For the window of $\alpha_+$, the exponent of $z_+$ runs over $|W|$ consecutive values, with $|W|$ of size $a/(kY)$
and $J$ of size $b'$. These sizes, and the choice of the depth $\lfloor YD_1\rfloor$, make the share of the
previous lemma at most $5/Y$.

% before: pd-kappa-defects-segment-mean-theorem-a-plus-window
Restoring the factor at $x_+$ can lose good vectors of the orthogonal family. Summed over the window of
$\alpha_+$, the loss is at most a share $5/Y$ of the size of the family.

% before: pd-kappa-defects-segment-mean-theorem-a-plus
The badness of a level depends on $\alpha_+$ only through the factor of the level pairing at $x_+$. Summed
over the window of $\alpha_+$, it is bounded by the badness of the polynomial part times the number of
monomials in the $\xi_c$, plus a share $5/Y$ of the dimension.

% before: pd-kappa-defects-segment-mean-theorem-a-top-model
The mean value estimate is proved for a model of a level: a chain of letters with ties, in which the degrees
of the tied covector letters vary over a progression. The model records only the degrees, the ties and the
conditions on the weights, and its numbers $\varepsilon_i$ measure the length of the progression against
the degrees below each tie.

% before: pd-kappa-defects-segment-mean-theorem-a-top-mean
Over a progression of degrees, the badness of the polynomial parts is small on average. It is bounded by
the numbers $\varepsilon_i$ times the dimensions of the polynomial parts with one tied degree raised to the
top of the progression.

% before: pd-kappa-defects-segment-mean-theorem-a-top-box
The tuples compatible with a digit class form a box, one interval of a congruence class for each entry.
Moving an entry within its set keeps the block and the level, so the level condition holds along the box.

% before: pd-kappa-defects-segment-mean-theorem-a-top-attached
A digit class, a tuple, a block and a level give a model instance. Its letters are those of the block with
the lowered degrees, its ties are the tied pairs $x_{u_i},y_{u_i}$, and its progressions are the possible
degrees of the $y_{u_i}$ as the entries $\alpha_{u_i}$ run over the box.

% before: pd-kappa-defects-segment-mean-theorem-a-top-sizes
The hypotheses of the mean value estimate include bounds on the degrees of the attached model instance and
on the sums $S_i$ of the degrees below each tie. They come from the bounds on the block and on the window.

% before: pd-kappa-defects-segment-mean-theorem-a-top-hyp
With these bounds, the attached model instance satisfies all the model hypotheses. The remaining ones
concern the arrangement of the letters and the ties, which is that of every block.

% before: pd-kappa-defects-segment-mean-theorem-a-top-eps
The mean value estimate is applied with the factors $\varepsilon_i$ of the attached model instance, which
must be of size $1/Y$. As $\varepsilon_i$ is $2S_i/R_i$ plus a term divided by $E=Yc_S$, this needs the
progressions to be long compared with the sums $S_i$.

% before: pd-kappa-defects-segment-mean-theorem-a-top-fibre
Fixing the entries of $\alpha$ outside the tied positions leaves a fibre of the box, which the degrees of
the $y_{u_i}$ identify with the progression of the attached model instance. On each fibre, the mean value
estimate for the model is the estimate for the level.

% before: pd-kappa-defects-segment-mean-theorem-a-top
This is the mean value estimate for the polynomial parts of a level over the tuples compatible with a
digit class. The badness of the polynomial part is bounded by $(30+C_M)/Y$ times the dimensions of the
polynomial parts at the tuples with one entry raised to the top of the window.

% before: pd-kappa-defects-segment-mean-theorem-a-directions-independent
The generators of different levels of a block are told apart by their exponents at the variables $z'_l$
and $z'_+$. A generator of level $(j,j_+)$ has these exponents at most $(j,j_+)$, with equality in a term
from which the generator can be recovered. So the spans of the levels form a direct sum.

% before: pd-kappa-defects-segment-mean-theorem-a-directions
The dimensions of the polynomial parts, weighted by the numbers of monomials in the $\xi_c$, add up over
the levels of a block to at most the dimension of the block. With condition (11), this bounds the total
weight of a digit class by $d(\delta;\alpha)$.

% before: pd-kappa-defects-segment-mean-theorem-a-segment
The estimates on the levels combine into one bound for each digit class: the badness of its blocks,
summed over the window, is a fraction $\varepsilon$ of the sum of the numbers $d(\delta;\alpha)$. The
fraction is of size $(t+1)/Y$, and the badness also bounds the rank lost in passing to good sublattices of
depth $D_A$.

% before: pd-kappa-defects-segment-mean-theorem-a
From condition (11), the mean of the segment defect over the window is of size
$(S_j-\eta_j+1)N_1/Y$; the statement also assumes conditions (1) to (10), which its proof does not use. The constant $C_A$ may depend on $C_\kappa$ and $C_s$, and the depth $D_A$ of the
good shares is formed with the same constant.

% before: pd-kappa-defects-segment-mean-letters-families
Conditions (1) and (3) concern projectors attached to representations of $\mathfrak{gl}_n$ on spaces of
polynomials. Their proofs use only the commutation relations, local finiteness and the weights, so they
are carried out for an abstract family of operators with these properties. Homomorphisms between two
such families, and their invariants, enter through the integrality bound.

% before: pd-kappa-defects-segment-mean-letters-span
Every vector of such a family is a sum of vectors generated by highest weight vectors. This is Weyl's
theorem in the form used here, and it determines an isotypic projector by its values on the isotypic
subspaces.

% before: pd-kappa-defects-segment-mean-letters-forms
The Fock pairing is contravariant for the letter operators. A contravariant form separates the isotypic
subspaces, so a nondegenerate one makes their sum direct and the isotypic projectors self-adjoint.

% before: pd-kappa-defects-segment-mean-letters-trace
The integrality of an isotypic projector is reduced to that of an invariant projector, which is bounded by
Lemma~\ref{pd-gain-weyl-invariant-denominator}. The reduction writes the isotypic projector through the
invariant projector of a space of homomorphisms, using Schur's lemma on the irreducible representation
generated by one highest weight vector.

% before: pd-kappa-defects-segment-mean-letters-separators
A lattice stable under the binomials in the diagonal operators contains the weight components of its
elements. The binomials in one diagonal operator separate any two weights by functions with integer values.

% before: pd-kappa-defects-segment-mean-letters-integral
The isotypic projectors of a family have bounded denominators on a lattice stable under the divided
powers. The exponent of the denominator depends only on $n$, on the spread of the weight and on the root
length of the representation.

% before: pd-kappa-defects-segment-mean-letters-operators
The letter operators of a lettering act on the polynomials in its rows of variables. They form a
$\mathfrak{gl}_n$-family for every set of letters, and the Fock pairing is contravariant and nondegenerate
for them.

% before: pd-kappa-defects-segment-mean-letters-lattices
Condition (1) uses two integral structures, the polynomials with integral coefficients and the module
$\mathcal B$, whose coefficient at $z^\nu$ may have the denominator $p^{v_p(\nu!)/2}$. Both are stable under
the divided powers of the letter operators, and both are finitely generated in each row degree.

% before: pd-kappa-defects-segment-mean-letters-polynomials
The last lemmas combine to bound the isotypic projectors of the letter operators on both integral
structures. The exponent depends on $n$ and on the total degree of the letters involved.

% before: pd-kappa-defects-segment-mean-letters-degrees
The hypotheses of an instance bound the degrees of the letters of a kept block by a multiple of $am$. This
bounds the exponent of Lemma~\ref{pd-kappa-defects-segment-mean-letters-polynomials} by $\kappa_1$.

% before: pd-kappa-defects-segment-mean-letters-highest
Condition (3) asks for a projector onto the vectors of one weight that are highest for the first $r-1$
indices. On the weight space it is an isotypic projector of the smaller family, and the contravariant
form identifies its kernel with the image of the lowering operators.

% before: pd-kappa-defects-segment-mean-letters
Conditions (1) and (3) concern only the letters of the kept blocks of the segment and the exponent
$\kappa_1$. They hold for every instance once the constants are large.

% before: pd-kappa-defects-segment-mean-weyl-lattices
A lattice stable under the divided powers that meets the highest weight space in $\mathcal Ov$ contains
every vector generated from $v$ and lies in $\nabla(\lambda,v)$. This is the second half of condition (2),
and it uses only the weights.

% before: pd-kappa-defects-segment-mean-weyl
Condition (2) places every integral structure on an irreducible module of $\mathfrak{gl}_h$ that meets
the highest weight space in $\mathcal Rv$ between $\Delta(v)$ and $\nabla(\lambda,v)$. The factor
$p^{\kappa_1}$ bounds the gap between these two lattices when the highest weight has spread at most $8a$.

% before: pd-kappa-defects-restriction-universal
Membership in $H_{r+1}(c,\kappa)$ is an identity over every algebra. Over an integral domain it suffices
to check it once, for a generic lower triangular matrix and a generic matrix, since every other pair is
a specialization of these.

% before: pd-kappa-defects-restriction-basic
The spaces $H_{r+1}(c,\kappa)$ are closed under sums and products, with the characters adding up. The
determinant and the corner coordinate $g_0$ are members, and changing the base ring keeps a function in
its space.

% before: pd-kappa-defects-restriction-divide
The corner coordinate $g_0$ transforms by the first diagonal entry of the lower triangular factor. So a
quotient by $g_0$ of a function in $H_{r+1}(c,\kappa)$ lies in the space with $c$ lowered by one.

% before: pd-kappa-defects-restriction-prime
Divisibility questions in the coordinate ring of $GL_n$ reduce to the polynomial ring, since inverting the
determinant or the corner entry $x_{11}$ does not affect a prime of the polynomial ring whose value at the
identity matrix is not a unit. Over a discrete valuation ring, a function whose reduction vanishes is divisible by the
uniformizer.

% before: pd-kappa-defects-restriction-polynomials
A polynomial member of $H_n(\nu)$, read through the transpose, is a polynomial in the dual Weyl module of
weight $\nu$. Conversely the transposed bideterminants of shape $\nu$ lie in $H_n(\nu)$, because a lower
triangular matrix acts on the first $s$ rows of a matrix through its upper left $s\times s$ block.

% before: pd-kappa-defects-restriction-standard
The dual Weyl module is spanned by bideterminants over every integral domain, with a basis indexed by
semistandard tableaux. The statement is needed for all $\nu\in\mathbf Z^n_{\geq0}$, since shifting a character
to make it nonnegative need not make it non-increasing.

% before: pd-kappa-defects-restriction-lift
A function of $H_{r+1}(c,\kappa)$ over the residue field lifts to the discrete valuation ring. After a power
of the determinant it is a combination of transposed bideterminants, and bideterminants have integer
coefficients.

% before: pd-kappa-defects-restriction-count
The dimensions of the spaces $H_{r+1}(c,\kappa)$ over a field are numbers of semistandard tableaux. Removing
the largest entry of a tableau gives the branching rule for these numbers, from which the effect of a shift
and of raising the first entry follow.

% before: pd-kappa-defects-restriction-dimension
Over a field, multiplying by a power of the determinant identifies $H_{r+1}(c,\kappa)$ with a transposed dual
Weyl module. Its dimension is therefore the number of semistandard tableaux of the shifted weight.

% before: pd-kappa-defects-restriction-cell
On the open set where the corner entry is invertible, every matrix factors as a lower triangular matrix
times a matrix whose first column is the first unit vector. So a function in $H_{r+2}(c,\kappa)$ is
determined by its restriction to the second kind of matrix, which is the evaluation $\varepsilon$ of
condition (7).

% before: pd-kappa-defects-segment-mean-restriction
Conditions (4) to (7) are statements about the functions on $GL_{r+1}$ that transform by a character,
with no reference to the instance. They describe division by the coordinate $g_0$ and the restriction to
the matrices whose first column is the first unit vector.

% before: pd-kappa-defects-paths-reducible
The operators on the letters satisfy the commutation relations of $\mathfrak{gl}_r$, and the polynomial
ring is the sum of its weight spaces. In characteristic $0$ a finite-dimensional representation of this
kind splits into isotypic components.

% before: pd-kappa-defects-paths-isotypic
The spaces $Y(\mathcal L,M)$ of condition (8) are defined as the largest stable subspaces with a weight
bound. They are the sums of the isotypic components whose highest weights have first entry at most $M$,
which is what allows a comparison with the path generators.

% before: pd-kappa-defects-paths-lower-types
Membership in a sum of isotypic subspaces is a condition on the projections of a polynomial. It can be read
off from the dominant weights of the subspace the polynomial generates, which is the form of the condition in
the comparison with the path generators.

% before: pd-kappa-defects-paths-lower
Over a field, the path generators with a constraint on one vector of the path span the polynomials whose
isotypic components for the corresponding initial letters have weights in a lower set. The integral
statement that follows is deduced from this one through the path basis.

% before: pd-kappa-defects-paths-nested
Over the discrete valuation ring, the integral polynomials whose isotypic components lie in prescribed
lower sets of weights are spanned by the path generators whose paths pass through those sets. The lower
sets are taken for the dominance order, and condition (8) uses the sets cut out by a bound on the first
entry.

% before: pd-kappa-defects-segment-mean-paths
Condition (8) describes the integral polynomials whose weights for initial sets of letters are bounded.
They are spanned by the path generators of Definition~\ref{pd-symbols-path-generators} whose paths obey
the same bounds.

% before: pd-kappa-defects-segment-mean-ratio
Condition (9) is the lower bound for the product of the ranks of the segment lattices together with the
bound on the margin term. Both hold for exponents in the enlarged windows, since $x_a$ and $x_b$ are at
most $2J_0$.

% before: pd-kappa-defects-digits-half
The lower bound for the kept blocks with given digits carries a factor that is close to $1$ when the margin
is small. Under the size relations of the hypotheses this factor is at least $1/2$ for every margin up to
$b'/(kQY)$.

% before: pd-kappa-defects-digits-chains
The rank of the kept blocks with given digits is computed by counting. Each level of a block contributes
a space of polynomials with weight bounds, whose dimension is a sum of Weyl dimensions over the paths that
respect the bounds, times the number of monomials in the variables $\xi_c$.

% before: pd-kappa-defects-digits-block-rank
The generators of a block are products of a polynomial in the letters, a polynomial in the $\xi_c$ and
fixed factors that record the level. The dimension of their span is the sum over the levels of the
products of the dimensions of the two factor spaces.

% before: pd-kappa-defects-digits-weyl
The dimension of a space spanned by path generators is the number of pairs of a path and a standard
filling. For a fixed path, the standard fillings are semistandard tableaux of the shape of its end, and
their number is the Weyl dimension of that end.

% before: pd-kappa-defects-digits-chain-dimension
The polynomials of given row degrees that satisfy weight bounds for initial sets of letters have a basis of
path generators. Counting the standard fillings of a path by the Weyl dimension of its end gives the
dimension as a chain sum.

% before: pd-kappa-defects-digits-count
Together, the two dimension formulas turn the rank of the kept blocks with given digits into the chain
count. The bounds of conditions (10) and (11) then become bounds for this count.

% before: pd-kappa-defects-digits-branching
The upper bound compares the chain count with a sum of Weyl dimensions over longer chains of one rank more.
Summing out the last vector of such a chain uses the branching of Weyl dimensions from rank $r+1$ to
rank $r$.

% before: pd-kappa-defects-digits-top-rows
After the letters $x_i$ with $i\leq\eta_j$, the chain conditions give no bound on the first entry of a
chain. The bound comes instead from the degrees $w_l$ of a kept block, whose partial sums stay below the
capacities.

% before: pd-kappa-defects-digits-count-upper
The chain count of the kept blocks with given digits is at most the product of the coarse dimensions. This
is a statement about paths and Weyl dimensions only, with no reference to the ring.

% before: pd-kappa-defects-digits-count-pairs
The proof of the upper bound for the chain count maps each source injectively to a target. For the lower bound the
targets are mapped back to sources, which is possible only for targets whose paths respect the bounds that define a
kept block. Those targets are built from the kept pairs.

% before: pd-kappa-defects-digits-count-binomial
Two bounds on binomial coefficients enter the final count. The first controls the loss when the degree $\hat a$ is
lowered by the margin, and the second the loss when only the levels with $j_+\leq\Lambda b_+/k$ are counted.

% before: pd-kappa-defects-digits-count-pair-number
The number of pairs is needed to bound the total weight of the pairs that are not kept. A path for the letters of a
segment starts at $0$ and passes few letters, so each of its vectors has few nonzero entries, and this bounds the number
of pairs.

% before: pd-kappa-defects-digits-count-pair-facts
At each letter a pair splits into a first entry and a tail. The tails form the path of a source, and the first entries
record the capacities of its block, so the conditions for a kept block become inequalities between first entries.

% before: pd-kappa-defects-digits-count-coarse
The coarse degrees $\hat b_i$ of a segment grow quickly with $i$, and they lie below or above $\hat a$ according as
$i\leq\eta$ or $i>\eta$. Both properties come from the gap between the odd powers of $p$ and $a/b'$ in hypothesis (2) of
Theorem~\ref{pd-gain-theorem-one}, and they decide which conditions for a kept block a pair has to meet.

% before: pd-kappa-defects-digits-count-pair-bound
In the setting of a segment the coarse degrees are at most $2a/p$. The count of the pairs then becomes a power of $8a$
with an exponent of order $N_1^2$.

% before: pd-kappa-defects-digits-count-weyl-bound
The Weyl dimension of a non-increasing vector is bounded by a product of binomial coefficients, one for each entry. The
bound comes from a single path that reaches the vector one entry at a time, whose weight is one term of a sum equal to
that product.

% before: pd-kappa-defects-digits-count-weight
In a pair, the last vector $\lambda$ has nonnegative entries at the start, zeros in the middle and nonpositive entries at
the end, and its first entries are bounded by the numbers $t_i$. With the bound for Weyl dimensions this bounds
$\operatorname{wd}(\lambda)$ by binomial coefficients in $\lambda_1$, the $t_i$ and the last entries of $\lambda$.

% before: pd-kappa-defects-digits-count-bracket-ratio
The coarse degrees $\hat b_i$ grow by a factor of about $p^2$ from one letter to the next, as the distances $d_{j,i}$ are
odd and increase. The coarse degrees $\hat\alpha_i$ of the covector letters grow in the same way, since all the $a_T(y_{j,i})$
lie in a narrow window above $a$. So their sums over the earlier letters are small multiples of the next one.

% before: pd-kappa-defects-digits-count-bracket-facts
For the bound on a pair that is not kept, the entries that enter the Weyl dimension are compared with the coarse degrees.
The first entries $c_l$ and $t_i$ and the last entries $n_i$ of a pair are nonnegative, and each is at most a coarse degree
up to the small sums of the earlier ones.

% before: pd-kappa-defects-digits-count-bracket-large
Each factor of the bound for a pair that is not kept is shifted by $m$, and the shift has to be absorbed into a factor
close to $1$. This needs the coarse degrees to be large compared with $k^3$ and $p$ to be large compared with $k^2$, which
the hypotheses on the constants give.

% before: pd-kappa-defects-digits-count-bracket-core
The bound for a pair that is not kept rests on an inequality between real numbers. The numbers $c_l$ change from one
index to the next by the coarse degree minus $\bar o_i$, and the failure of a condition for a kept pair at one index costs a
factor $1-\Lambda/(4k)$ or $1-\Lambda/k+1/k^2$ there. Raised to the power $m-1$, that factor gives $e^{-\Lambda/5}$, up to
a bounded loss from the shifts.

% before: pd-kappa-defects-digits-count-bracket
When a pair is not kept, one of the two conditions for a kept pair fails at some index. The failure forces the numbers
$t_i$ or the last entries of $\lambda$ to be smaller, by a factor depending on $\Lambda$, than the coarse degrees they are
compared with.

% before: pd-kappa-defects-digits-count-not-kept
A pair that is not kept carries a Weyl dimension that is a very small multiple of the product of the coarse dimensions.
Together with the count of the pairs, this makes the share of the pairs that are not kept in the total weight small.

% before: pd-kappa-defects-digits-count-tail
The total weight of the pairs is the product of the coarse dimensions, by the branching rule and the dimension count of
the chain sums. The pairs that are not kept are few and light, so the kept sum retains all but a fraction
$e^{-\Lambda/8}$ of the total.

% before: pd-kappa-defects-digits-count-kept
A kept pair, with a level $j_+$ for the last letter, gives a block, a level and a path, by reversing the map in the proof
of Lemma~\ref{pd-kappa-defects-digits-count-upper}. The conditions for a kept pair are what makes the block kept and the
path respect the thresholds of its level.

% before: pd-kappa-defects-digits-count-match
The map from kept pairs to sources is injective and preserves the weights. So the kept sum, times the sum of the
binomial coefficients of the counted levels $j_+$, is at most the chain count.

% before: pd-kappa-defects-digits-count-lower
When the margin is small compared with $\hat a_j$, the chain count is also bounded below by a fraction of
the product of the coarse dimensions. The fraction is the factor of
Lemma~\ref{pd-kappa-defects-windows-model-segment}.

% before: pd-kappa-defects-digits-upper
The kept blocks with given digits contribute at most the product of the coarse dimensions to the rank of
the segment lattice, for every margin $\mu\geq0$. It is the upper half of the two-sided bound from which
conditions (10) and (11) follow.

% before: pd-kappa-defects-digits-lower
Conversely the kept blocks with given digits contribute at least a fraction of the product of the coarse
dimensions, provided the margin is small compared with $\hat a_j$. The fraction is the factor of
Lemma~\ref{pd-kappa-defects-windows-model-segment}, and Lemma~\ref{pd-kappa-defects-digits-half} shows
that it is at least $1/2$ in the range of condition (11).

% before: pd-kappa-defects-segment-mean-digits
Conditions (10) and (11) bound the part of the rank of a segment lattice that comes from the kept
blocks with given digits. It lies between half of the product of the coarse dimensions and that product
when the margin is at most $b'/(kQY)$, and below that product when the margin is $0$.

% before: pd-kappa-defects-segment-mean
The segment defect depends on the exponents $a_T$ only through the values at the positions of its
segment. Its mean over a window for these values is of size $(S_j+1)N_1/Y$.

% before: pd-kappa-defects-numerical
The lemmas above are stated for an integer $a$ and with the constant $C_H$ of their own thresholds. The
data of Theorem~\ref{pd-gain-theorem-one} satisfy their hypotheses with $\bar a$ and $Y_b/2$ once $C_H$
is large, and the enlarged windows $I'_a$ and $I'_b$ lie in the windows those lemmas allow.

% before: pd-kappa-defects-margin-share
The margin term is bounded for exponents in the windows of the rank bound and for every margin $\mu$
between $0$ and $b'/(kQY)$. The margin enters the term only through $kQ\mu/b'$, which is then at most
$1/Y$, so the bound at $m_0$ gives this bound with the constant increased by $C_\kappa$.

% before: pd-kappa-product-defects
The defect on the product spaces is the sum of a margin term and of one defect for each segment
between consecutive even elements of $T$. A segment defect depends on $A$ only through its values at
the $S_j+1$ positions of the segment, and its mean over these values is of size $(S_j+1)N_1/Y$.

% before: pd-kappa-product-segments
Summing a segment defect over $I_a'^T$ multiplies its sum over the values at the positions of the
segment by the number of choices at the other positions. So its mean over $I_a'^T$ is its mean over
the segment, of size $(S_j+1)N_1/Y$.

% before: pd-kappa-product-rank
The mean of the flat defects is controlled through the product spaces. On them the rank at an edge
is bounded below with the extremal projectors, by a defect whose mean over slightly enlarged windows
is of size $|T|N_1/Y$.

% before: pd-kappa-numerical
For a large constant $C_H$, the hypotheses of the theorem imply a list of inequalities among the
numbers of the construction. They are the properties of the data used in the two propositions that
follow.

% before: pd-kappa-projection
The projection property for the function $\kappa$ follows from the definition of the flat defects
and the passage from flat weights to the given weights. When $\kappa$ is at least $1$ it holds with
no edges.

% before: pd-kappa-mean
The mean bound compares the flat defects of the pieces of a product space with the defect on the
product space. Averaging over the windows, and using that most exponents stay in the windows after a
shift by $c_{cl}$, gives a bound of size $|T|N_1/Y$.

% before: pd-gain-theorem-one
For primes between $p_0$ and $X'$ the projection property is proved with edges of index $p^{n'}$,
for exponents in short windows. The defect $\kappa$ now depends on the exponents, and only its mean
over the windows is bounded.

% before: pd-gain-coherence
The local models at the places of the gain data are built from the global unitary group of the
forms in degree $0$. Group coherence records how: the components are classes of a level of that
group, the groups of the components are its stabilizers, and the comparison of the forms with the
local functions has coefficients in the coefficient field.

% before: pd-gain-projection-monotone
The two gain theorems give the projection property with values of $\kappa$ that differ from the
local constants of the gain data. A larger $\kappa$ weakens the property, and so does replacing
$\kappa$ by its minimum with $1$.

% before: pd-gain-good-segments
Theorem~\ref{pd-gain-theorem-one} needs the even top positions of an orbit to be spread out, with
few odd positions between consecutive ones. For a tiled set of positions this follows from the
marked tiles and the window bound of its definition.

% before: pd-gain-second-hypotheses
At primes below $p_0$ the gain data use the projection property with explicit thresholds built from
the scale and the constants $K_G$. Theorem~\ref{pd-gain-theorem-two} gives it once its parameters are
matched with those of the scale, and the matching is a comparison of the constants with $C_0$.

% before: pd-gain-first-hypotheses
At primes between $p_0$ and $X'$ the local constants are defined as the least values for which the
projection property holds. Their mean is then at most that of the function of
Theorem~\ref{pd-gain-theorem-one}, once the hypotheses of that theorem are checked for the windows of
the gain data.

% before: pd-gain-weyl-comparison
The integral module of a local model is built from Weyl modules and the projection property from
dual Weyl modules. Over the ring of a Hida setting the two differ by a bounded power of $p$.

% before: pd-gain-weyl-sandwich
Summed over the positions of a local model, the comparison of Weyl modules bounds the power of $p$
that separates the integral module from its dual-Weyl version. The Weyl exponents of the gain data
are chosen at least this large and of size $fk^2\log_pa$.

% before: pd-gain-quantities
The law of the gain data is stated in terms of numbers computed from the gain data alone: the gains
at each prime, the classes counted there, the depth $\tau$ and the index excess $\iota$. The size data
of a family take their gains, depths and excesses from these numbers.

% before: pd-gain-coherence-consequences
Three consequences of group coherence enter the local theory. Trivial stabilizers enter the gluing
of the served lattices, the standard rings allow the gain theorems to be applied to the local models,
and rationality lets the served local lattices be spanned by images of global forms.

% before: pd-gain-base-change
The gain theorems are proved over the ring of an unramified quadratic setting, while the local
models live over a finite extension of it. The projection property survives this extension because
every lattice in it is formed by operations that commute with flat base change.

% before: pd-gain-admissibility
The law of the gain data applies to families that are admissible. The conditions bound the orbits
and the bands at each prime, ask for the projection property at the served orbits and for the Weyl
sandwich, and control the served local lattices.

% before: pd-gain-orbit-bounds
The law needs the projection property at each served orbit of each local model. It comes from the
two families of bounds, applied to the unramified quadratic setting under the orbit and carried to
the local model by base change.

% before: pd-gain-holds-high-classes
Above $X''$ the exponents $a'_s$ are too small, relative to the fills, to contribute to the class
sums at more than one power of $p$. Only the primes just above $X''$ contribute at all, and each
contributes a small amount.

% before: pd-gain-holds-third-numerics
The sums over primes in part (6) of the law are bounded by the number of primes in a range up to
about $X''$, times a bound for each prime. At large scales the resulting explicit functions of the
scale are at most $a\,x^{-16}$ and $b\,x^{-100}$.

% before: pd-gain-holds-fourth-numerics
Part (4) of the law sums, over the primes, the excess of the class sums over the gains. Each prime
contributes a few explicit terms, and their sums over ranges of primes are compared with
$a/x$ at large scales.

% before: pd-gain-law
For an admissible family, the law of the gain data bounds the class sums at every inert prime by the
gains, up to the lower terms. It also bounds the mean local constants at the candidate primes, and
compares the two versions of the class sums.

% before: pd-gain-weyl-column-projector
The comparison of Weyl modules rests on a projector onto the invariants of a representation of
$\mathfrak{gl}_n$ whose denominators are powers of $p$ of controlled size. It is built one column at
a time, from sums of composites of a raising and a lowering operator, whose coefficients are
inverses of binomial coefficients.

% before: pd-gain-weyl-invariant-denominator
Applying the column projectors one after another to the weight-zero part of a vector gives its
invariant component. The denominators add up to $p^{N_0d}$, with $N_0=n(n-1)/2$ the number of
pairs of columns.

% before: pd-gain-weyl-highest-lattices
For an irreducible representation, the smallest and the largest lattices with a given integral
highest weight line are compared through the endomorphism ring. The invariant projection of a
rank-one map onto the highest weight line is a scalar, and its denominator bound gives the
comparison.

% before: pd-gain-link-surplus
The surplus of a served representative counts what the unit gain exceeds over its Hecke loss and the
local gain. Its bound by $2mN(p,t_0)a/p$, condition (2) of the gain link, needs only the sizes of the
exponents and offsets and two of the parameter choices.

% before: pd-gain-link-admissible-places
The local models are attached only to the places admissible for a word. A prime with served
representatives is a candidate prime, and the sizes of the host put every place over it in this
class.

% before: pd-gain-link-descent
The comparison of a lattice instance carries vectors of the coefficient field to functions on the
local model. Integrality and lengths of lattices can be read on either side.

% before: pd-gain-link-served-lattices
Condition (1) of the gain link compares the lattice $N^\circ_u$ chosen with the size data with the
served local lattice of the local model. The comparison and its descent turn the bounds of the
admissibility of the gain data into an inclusion and a length bound at the place.

% before: pd-gain-holds-series
The class sums and the local gains are sums, over odd exponents $e$, of positive parts of $y/p^e$
minus a weight. Comparing them needs the geometric bound for such a sum and the effect of shifting the
weights or the numerator, which is felt only at the $e$ with $p^e<A$.

% before: pd-gain-holds-numerics
The mean bounds of the law need the scale to be large compared with the constants of the gain data.
The inequalities used compare powers of $x$ with powers of $\log x$ and of $x^\varepsilon$, and they
follow from the size bounds of the family and the bounds on $\Lambda$ and $Q$.

% before: pd-gain-holds-hypothesis-one
Condition (8) of admissibility bounds the local constants at a candidate prime $p\geq p_0$ only where
Hypothesis I holds. At large scales it holds at every good representative, by a computation from the
orbit sizes of condition (1), the band bounds of condition (4) and the sizes of $\Lambda$ and $Q$.

% before: pd-gain-holds-star-numerics
The comparison at one good representative ends in an inequality between explicit functions of the
scale. At large scales the quantities $b'$ and $\beta$ grow faster than every other term in it,
and the constant $C$ only has to dominate a few fixed combinations of the constants of $K_G$.

% before: pd-gain-holds-star
The inequality of the law at a candidate prime is assembled from one inequality for each good
representative. It compares the class sums of the top classes with that representative with the unit
gain, after the Hecke loss and the thresholds $t(p,t_0)$ and $D(p)$ are paid.

% before: pd-gain-holds-gain-bound
Summing the inequalities at the good representatives gives the inequality of the law at the
candidate primes, with the weights $\mathrm{fill}_j+\kappa_j$. At the other inert primes the local
gains bound the class sums, since the exponents excluded at a representative contribute nothing to
the class sums there.

% before: pd-gain-holds-transfer
The length bound of condition (11) is proved by covering a power of $p$ times a lattice by a
submodule and a few further elements. Such a covering passes to a smaller lattice, at the cost of the
exponent of a sandwich between the two.

% before: pd-gain-holds-length
Coverings at a sequence of exponents bound the length of a quotient. When the number of further
elements needed at the exponent $\tau+\beta S$ decays geometrically in $\beta$, their total
contribution is at most twice the first one.

% before: pd-gain-holds-esymm
The number of further elements in the coverings of condition (11) is an elementary symmetric function
of the local constants. The following inequality compares it with a power of their sum, the form in
which the deviation $\Sigma(p)$ enters the length bound.

% before: pd-gain-holds-walk
The covering of $M^\nabla$ by products of Hecke operators comes from a walk through a system of
submodules indexed by labellings of the served orbits. At each orbit the submodule can be enlarged to
a source module or shrunk to a target module, and the Hecke operator at that orbit carries the first
into the second. The walk moves a set of $\beta+1$ orbits upwards and applies each Hecke operator up
to $m$ times, so further elements are needed only at $m^{\beta+1}$ places for each such set.

% before: pd-gain-holds-count
The covering at one orbit is read off from the elementary divisors of the span of the vectors of the
projection property. A matrix identity with $q$ on the diagonal forces many of these elementary
divisors to divide $q$.

% before: pd-gain-holds-smith
Applied to a lattice between $\varpi H$ and $H$, the count gives a basis of the lattice all but a few
of whose members, multiplied by $q$, lie in the span of the given vectors. The number of exceptions is
at most the dimension minus the rank $n$ of the pairing.

% before: pd-gain-holds-local-chain
At one orbit the three labels stand for the lattice $H(L)$, the module $N_{\mathrm{src}}(L)$ and the
module $P_{\mathrm{tgt}}(L)$ of vectors that pair integrally with $N_{\mathrm{tgt}}(L)$. They form a
chain whose steps cost $\varpi_0^t$, an edge carries the source module of its source into the target
module of its target after multiplication by $p^A$, and the group action moves all three with the
lattice.

% before: pd-gain-holds-systems
Choosing one of the three modules at each orbit and imposing the conditions of the integral module
gives a module of functions for each labelling. The functions supported on one orbit of lattice
tuples are built from a single value by the group action.

% before: pd-gain-holds-orbits
With trivial stabilizers and finitely many classes of lattice tuples, a function in one of these
modules is determined by its values at a set of representatives. This bounds the number of
representatives by the rank of $M^\nabla$, and reduces the covering to one lattice tuple at a time.

% before: pd-gain-holds-ball-dual
The Hecke operator at an orbit carries the source module there into the target module, since $p^A$
times an element of a source module pairs integrally with the target modules of the edges above it.
Applied to a function concentrated on one orbit of lattice tuples, it reproduces the value at the
chosen edge, and the other edges contribute multiples of $p^E$ when the value is divisible there.

% before: pd-gain-holds-local-smith
The projection property at an orbit supplies the vectors to which the Hecke operator is applied.
Through the elementary divisor count, all but $\kappa\dim V$ members of a basis of the target module
become, after multiplication by $p^D$, combinations of these vectors.

% before: pd-gain-holds-covering
At a set $B$ of served orbits, a product of local basis vectors that avoids the exceptional ones at
some orbit of $B$ comes from a Hecke operator. The other products number at most
$\prod_{v\in B}\kappa_v$ times the dimension, and summing over the representatives and applying
Nakayama's lemma gives the covering of the module $M_t(B)$.

% before: pd-gain-holds-selection
The served local lattice is spanned by the images of $M$ under products of the Hecke operators at the
served orbits, multiplied by powers of $p$. The covering of $M^\nabla$ by these products needs a
number of further elements bounded by the local constants of the projection property at the served
orbits.

% before: pd-gain-holds-hecke-rows
An edge at an orbit changes the lattice there by a diagonal matrix in a suitable basis. The effect on
the Weyl and dual Weyl modules is read off from the degrees of their elements in the rows of $x$, and
these degrees are bounded by the partial sums of $\nu$.

% before: pd-gain-holds-hecke-step
Dividing the first $r$ vectors of a basis by $c$ enlarges the lattice, and the lattices of polynomials
attached to it change with it. By the row bound of Lemma~\ref{pd-gain-holds-hecke-rows}, the factor
$c^{\lambda_1+\cdots+\lambda_r}$ carries the old lattice of polynomials into the new one.

% before: pd-gain-holds-hecke-stable
The products of Hecke operators preserve $M$ and $M^\nabla$ when each Hecke operator does. The Hecke
operator at an orbit $v$ sums the values at the edges below $L_v$, multiplied by $p^{A_v}$.

% before: pd-gain-holds-served-length
At a place with served representatives, the length bound of condition (11) follows from the covering.
The covering of $M^\nabla$ passes to $M$ through the Weyl sandwich, the deviation $\Sigma(p)$ counts
the further elements, and the bound on $M$ restricts to the ramified part because it is saturated.

% before: pd-gain-holds-window-sizes
Condition (5) of admissibility bounds $\Lambda$ and $\log Q$ by powers of $x$. Both are controlled by
the length $G'$ of the windows and the count $N_1$ of Definition~\ref{hf-good-orbit}, which are of
size $x^{\varepsilon/8}\log x$ and $\log x$.

% before: pd-gain-holds-condition-five
With these sizes, condition (5) holds as soon as a fixed multiple of $(\log x)^3$ is at most
$x^\delta$. The rank $k$ enters only through $p_0=k^2x^{26}$, which a polynomial bound on $k$ keeps
below a power of $x$.

% before: pd-gain-holds-hypothesis-two
At the candidate primes below $p_0$, the projection property of the orbits comes from
Lemma~\ref{pd-gain-second-hypotheses}, which has a long list of numerical hypotheses. At large scales
they follow from the size bounds and from the lower bound on the orbit sizes in condition (1).

% before: pd-gain-holds-band-sums
Condition (4) asks that few primes have a power close to $X'$. For a fixed exponent such primes lie
in a short interval, and the exponent is at most about $2\log X'$, so the sum of $\log p/p$ over them
is small when the band is narrow.

% before: pd-kappa-orbit-large-scale
Hypothesis (I) of Lemma~\ref{pd-gain-orbit-bounds} asks that the windows of the tiled sets be shorter
than the tiles and that the orbits be long compared with $G'$. Both hold at large scales: a prime of
inert class exceeds a fixed power of $x$, by the splitting bound of the host, so the window length
$L_2$ stays bounded, and $G'$ grows only like $x^{\varepsilon/8}\log x$.

% before: pd-gain-holds-positions
Condition (1) of admissibility concerns the positions of the top classes along one orbit, which form
a tiled set of positions when the representative is good. Marked tiles occur among any $M_0+1$
consecutive tiles, and each of them holds even elements of the set, so even elements recur within
distance $G'$. The bounds on the windows limit how many elements lie close together.

% before: pd-gain-holds-orbit-sizes
Counting a tiled set by its windows bounds its size from both sides, in terms of $G'$ and $N_1$. These
two numbers are tied to $\log x$ by their definitions: $G'\delta_1$ is at most a constant times
$\log x$, and $N_1$ is at most a fixed multiple of $G'\delta_1$.

% before: pd-gain-holds-orbit-condition
At a prime of inert class the top classes with a given representative $t_0$ are the elements
$t_0+ih_p$ of its orbit that lie in $T$, indexed by the positions. Read along the orbit, the
statements about a tiled set become the three parts of condition (1).

% before: pd-gain-holds-two-adic
Condition (3) concerns the primes of inert class whose element $h_p$ has small order. The image of
$h_p$ in $H_e$ is the class of $p$, and an odd integer whose class in $H_e$ has small order is
congruent to $\pm1$ modulo a large power of $2$.

% before: pd-gain-holds-order-sums
Primes congruent to $\pm1$ modulo a large integer $m$ are sparse: the sum of $\log\ell/\ell$ over
those up to $Y$ is of size $(\log Y)^2/m$. Taking $m$ as large as the bound $J$ on the order allows
gives a bound of size $J(\log Y)^2/2^e$.

% before: pd-gain-holds-small-orders
At a scale $x$ the order $2^e$ of $H_e$ is nearly $x^{2+8\varepsilon}$, while $\log X''$ is about
$x^{\varepsilon/3}$. With $J=x^{\varepsilon/4}$ the bound of Lemma~\ref{pd-gain-holds-order-sums} is
then at most $x^{-2-7\varepsilon}$ once $x$ is large, which is condition (3).

% before: pd-gain-link-holds
The gain link is verified for the designed families of Chapter~\ref{chapter-assembly}. Its proof
reduces to the local theory of Frobenius orbits at each prime, through the coherence of the group
structures of the lattices and a count of the top classes whose orbits are not good at some prime.
